% ============================================================
% IDEIGLENES főfájl — amíg a ftsrg/thesis-template-latex sablon
% be nem kerül a projektbe. Akkor ezt a fájlt a sablon thesis.tex-e
% váltja ki, a chapters/ fájlok pedig \include-dal illeszthetők át.
% Fordítás: pdfLaTeX (Overleaf alapértelmezés).
% ============================================================
\documentclass[12pt,a4paper,oneside]{report}

\usepackage[T1]{fontenc}
\usepackage[utf8]{inputenc}
\usepackage{lmodern}
\usepackage[magyar]{babel}

\usepackage{graphicx}
\usepackage{xcolor}
\usepackage{listings}
\usepackage{enumitem}
\usepackage[hidelinks]{hyperref}

% --- Kódrészletek alapbeállítása ---
\lstset{
	basicstyle=\ttfamily\small,
	keywordstyle=\bfseries\color{blue!60!black},
	commentstyle=\itshape\color{green!40!black},
	stringstyle=\color{orange!70!black},
	numbers=left,
	numberstyle=\tiny\color{gray},
	frame=single,
	breaklines=true,
	showstringspaces=false,
	tabsize=4,
	captionpos=b
}
\renewcommand{\lstlistingname}{kódrészlet}

% --- Címlap-adatok (PLACEHOLDER — később kitöltendő) ---
\title{HUNting -- Magyar Vadászati Közösségi és Hirdetési Portál fejlesztése}
\author{SZERZŐ NEVE (placeholder)}
\date{2026}

\begin{document}

\maketitle
\tableofcontents

% --- Minden fejezet elkészült; a számozás természetes. ---

% !TeX root = ../main.tex
% 1. fejezet — Bevezetés
% A ftsrg/thesis-template-latex sablonba \include-dal illesztendő tartalomfájl.
% Hivatkozott források: Backend/HUNting.API/Models/Users/AccountType.cs,
%   thesis/planning/feature_matrix.md (funkciókör), Frontend/HUNting/src/routes/appRouter.tsx
% MEGJEGYZÉS: az 1.4 alfejezet konkurens megoldásainak konkrét megnevezése a szerzővel
%   egyeztetendő (ld. TODO-kommentek) — jelenleg kategóriaszinten, óvatos állításokkal íródott.

\chapter{Bevezetés}
\label{chap:bevezetes}

\section{Motiváció és problémafelvetés}
\label{sec:bev-motivacio}

A vadászat Magyarországon szervezett keretek között, jellemzően vadásztársaságokban folyó
tevékenység, amelynek mindennapjait sokrétű szervezési és adminisztrációs feladat kíséri. A
társaságoknak kommunikálniuk kell tagjaikkal a közös vadászatokról, munkákról és hírekről;
kezelniük kell a tagfelvételt; a vadászok eszközöket -- fegyvert, optikát, ruházatot,
felszerelést -- adnak-vesznek egymás között; a területi tájékozódáshoz megfigyelőpontok, etetők
és területhatárok ismerete szükséges; emellett a vadász saját vadászati eseményeiről is
nyilvántartást kíván (vagy köteles) vezetni.

Ezek az igények ma jellemzően egymástól független, az adott célra csak részben alkalmas
csatornákon oszlanak szét: a társasági kommunikáció telefonon, levelezőlistákon vagy közösségi
médiás csoportokban zajlik; az adásvétel általános célú hirdetési oldalakon vagy zárt
csoportokban történik, ahol a vadászati eszközök sajátos kategóriái nem jelennek meg; a területi
tudás gyakran papírtérképeken vagy egyszerűen a tagok fejében él; a naplóvezetés pedig
papíralapú vagy táblázatkezelős. A széttagoltság többszörös adatrögzítést, információvesztést
és belépési küszöböt jelent -- különösen az új tagok számára, akik a közösség hallgatólagos
tudásához nehezen férnek hozzá.

A dolgozat tárgya egy olyan webalkalmazás -- a \emph{HUNting} vadászati közösségi és hirdetési
portál -- megtervezése és megvalósítása, amely ezeket az igényeket egyetlen, integrált
rendszerben szolgálja ki: közös felhasználókezeléssel, egymásra épülő funkciókkal és egységes,
magyar nyelvű felülettel.

\section{A vadászati terület digitalizációs lehetőségei}
\label{sec:bev-digitalizacio}

A vadásztársadalom digitális eszközhasználatát kettősség jellemzi. Egyfelől a terület
hagyományőrző: a közösségi élet személyes jelenlétre épül, az adminisztráció sok helyen ma is
papíralapú. Másfelől az általános célú digitális csatornák -- közösségi médiás csoportok,
üzenetküldő alkalmazások, apróhirdetési portálok -- használata a vadászok körében is
mindennapos, ami jelzi, hogy a digitális eszközökkel szembeni ellenállás nem elvi, hanem
kínálati kérdés: az általános célú eszközök nem képezik le a terület sajátosságait.

Egy szakterületre szabott portál hozzáadott értéke éppen ebben rejlik. A tartományspecifikus
fogalmak -- vadásztársaság, tagsági kérelem, társasági szerep, vadfaj, megfigyelőpont, etető --
első osztályú entitásként jelenhetnek meg benne, nem pedig szabad szöveges bejegyzésekként egy
általános platformon. Az integráció további előnye az adatok összekapcsolhatósága: ugyanaz a
fiók ad fel hirdetést, igényel tagságot és vezet naplót; a társasági tagság jogosultságot
adhat közös térképi rétegekhez; a hirdetés tulajdonosa lehet maga a társaság is. Ezek a
kapcsolatok a széttagolt megoldásokban nem, vagy csak kézi egyeztetéssel valósulnak meg.

\section{Az alkalmazás céljai és célcsoportjai}
\label{sec:bev-celok}

A HUNting két célcsoportot szolgál ki, amelyeket a rendszer két, jogosultságaiban és
felületeiben eltérő fióktípusként kezel (a backend \texttt{AccountType} felsorolt típusa
szerint): az \emph{egyéni vadászokat} (Person) és a \emph{vadásztársaságokat} (Society).
E kettősség a rendszer egészén végigvonul, a regisztrációtól a menüszerkezetig
(\ref{chap:felhasznalok}.~fejezet). A funkcionális célokat hat funkcióterület valósítja meg,
amelyek egyben a dolgozat gerincfejezeteit is kijelölik:

\begin{enumerate}[nosep]
	\item \emph{felhasználókezelés, jogosultságok és adminisztráció} -- regisztráció a két
	      fióktípussal, bejelentkezés, profil, admin-felület (\ref{chap:felhasznalok}.~fejezet);
	\item \emph{hirdetési piactér} -- vadászati eszközök adásvétele kategóriákkal, szűréssel,
	      képfeltöltéssel (\ref{chap:hirdetesek}.~fejezet);
	\item \emph{vadásztársaságok és tagságkezelés} -- társaságkatalógus, tagfelvételi
	      kérelmek életciklusa, tag- és bejegyzéskezelés (\ref{chap:tarsasagok}.~fejezet);
	\item \emph{vadásznapló} -- személyes vadászati események rögzítése képekkel
	      (\ref{chap:naplo}.~fejezet);
	\item \emph{térképes funkciók} -- megfigyelőpontok, etetők, veszélyhelyek és területek
	      személyes és társasági rétegeken (\ref{chap:terkep}.~fejezet);
	\item \emph{AI-chatbot és időjárás} -- vadászati asszisztens és időjárás-tájékoztatás
	      (\ref{chap:chatbot}.~fejezet).
\end{enumerate}

Nem funkcionális célként a rendszer magyar nyelvű, reszponzív (asztali és mobil nézetben
egyaránt használható -- \ref{fig:home-responsive}.~ábra a~\ref{chap:alapok}.~fejezetben),
valamint szerveroldalon kikényszerített jogosultsági modellre épül. Fontos a célok őszinte
lehatárolása is: a rendszer demonstrációs érettségű -- éles üzemeltetésre, terheléses
méretezésre és teljes biztonsági keményítésre nem készült fel; az ismert korlátokat minden
funkciófejezet, összegzésüket a \ref{chap:teszteles}.~és a \ref{chap:osszefoglalas}.~fejezet
tartalmazza.

\section{Hasonló megoldások}
\label{sec:bev-hasonlo}

% TODO: a konkrét konkurens megoldások (nemzetközi vadásznapló-appok, hazai hirdetési
% portálok) megnevezése és jellemzése a szerzővel egyeztetendő; az alábbi szöveg
% kategóriaszinten, termékmegnevezés nélkül írja le a mezőnyt. Konkrét termékekről csak
% publikus, ellenőrizhető információ alapján szabad állítást tenni.

A HUNting által lefedett igényekre léteznek részmegoldások, ezek azonban jellemzően a terület
egy-egy szeletét célozzák. A \emph{vadásznapló-alkalmazások} (elsősorban nemzetközi mobil
appok) a vadászati események rögzítésére és részben térképes megjelenítésére szolgálnak,
közösségi és hirdetési funkciók nélkül, és magyar lokalizációjuk többnyire hiányzik. Az
\emph{általános hirdetési portálok} a vadászati eszközök adásvételét is hordozzák, de a
tartományspecifikus kategorizálás és a vadászközösségi kontextus nélkül; a fegyverhirdetések
kezelése ráadásul speciális szabályozási környezetet igényel. A \emph{közösségi médiás
csoportok} a társasági kommunikáció leggyakoribb mai csatornái: belépési küszöbük alacsony,
cserébe strukturálatlanok -- tagfelvételi folyamat, szerepek, naplózás vagy térképi rétegek
nem fejezhetők ki bennük natívan.

A HUNting elhatárolása e mezőnytől az \emph{integráció}: nem egy-egy részfunkcióban kíván a
specializált megoldásoknál többet nyújtani, hanem a vadászati közösségi élet funkcióit kapcsolja
össze egyetlen, közös fiók- és jogosultsági modellre épülő, magyar nyelvű rendszerben. A
választott felépítés sajátossága a kettős fióktípus: a társaság nem csupán címke vagy csoport,
hanem saját fiókkal, felülettel és jogosultságokkal rendelkező szereplő.

\section{A dolgozat felépítése}
\label{sec:bev-felepites}

A dolgozat szerkezete a rendszer funkcióközpontú jellegét követi. A keretfejezetek közül
a~\ref{chap:technologiak}.~fejezet a felhasznált technológiákat és választásuk indokait,
a~\ref{chap:architektura}.~fejezet a háromrétegű architektúrát és a rétegek közötti
kommunikációt, a~\ref{chap:alapok}.~fejezet pedig a minden funkció által közösen használt
infrastruktúrát -- adatbázis-kontextust, backend-csővezetéket, kliensvázat -- mutatja be.

A dolgozat gerincét a \ref{chap:felhasznalok}--\ref{chap:chatbot}.~funkciófejezetek adják.
Mindegyikük azonos belső sablont követ: a funkció céljának és felhasználói működésének
bemutatása után az adatmodellt, a backend-, majd a frontend-megvalósítást tárgyalja, végül a
teljes adatfolyamot, a biztonsági megfontolásokat és -- őszinte felmérésként -- az ismert
korlátokat összegzi. Ez a felépítés teszi lehetővé, hogy minden funkció a teljes vertikumában,
az adatbázistól a felhasználói felületig egy helyen legyen olvasható.

A gerincet két lezáró fejezet követi: a~\ref{chap:teszteles}.~fejezet a háromszintű
tesztkészletet és a futtatási környezetet dokumentálja, a~\ref{chap:osszefoglalas}.~fejezet
pedig az eredményeket összegzi és a továbbfejlesztési irányokat jelöli ki.

% !TeX root = ../main.tex
% 2. fejezet — Felhasznált technológiák
% A ftsrg/thesis-template-latex sablonba \include-dal illesztendő tartalomfájl.
% Hivatkozott források: Frontend/HUNting/package.json, tsconfig.app.json, vite.config.ts,
%   src/lib/queryClient.ts, src/App.tsx, src/index.css,
%   Backend/HUNting.API/HUNting.API.csproj, Program.cs, Data/HuntingContext.cs,
%   appsettings.json, src/pages/Map/MapPage.tsx, Services/OpenAiService.cs (hivatkozás),
%   src/services/weatherService.ts

\chapter{Felhasznált technológiák}
\label{chap:technologiak}

Ez a fejezet a rendszerben ténylegesen használt technológiákat mutatja be, rétegek szerint
haladva, minden választásnál megadva annak indokát is. A cél nem a technológiák általános
ismertetése, hanem annak rögzítése, hogy az adott eszköz \emph{ebben a projektben} milyen
szerepet tölt be; a funkcióspecifikus felhasználási részletek a gerincfejezetekben
(\ref{chap:felhasznalok}--\ref{chap:chatbot}.~fejezet) találhatók. A verziószámok a
csomagdefiníciókból (\texttt{package.json}, \texttt{HUNting.API.csproj}) ellenőrzött értékek;
a teljes technológiai palettát a~\ref{fig:stack}.~ábra tekinti át.

% TODO: T-01 ábra (figures/t-01-stack.pdf) — rétegzett stack-ábra: React/TS/Vite/Tailwind/
% React Query/Router/axios | ASP.NET Core 8/EF Core/Identity/JWT/Swagger | MSSQL | külső API-k.
\begin{figure}[htbp]
	\centering
	\fbox{\parbox[c][6.5cm][c]{0.9\textwidth}{\centering T-01: A technológiai stack áttekintése
	(elkészítendő)}}
	\caption{A HUNting technológiai rétegei}
	\label{fig:stack}
\end{figure}

\section{A frontend technológiái}
\label{sec:tech-frontend}

\subsection{React és TypeScript}

A felhasználói felület a React könyvtár 18-as főverzióján alapul~\cite{reactdocs}. A React
deklaratív modellje -- a felület a mindenkori állapot függvénye, a DOM-frissítéseket a könyvtár
vezérli -- jól illeszkedik egy adatvezérelt alkalmazáshoz, ahol ugyanazok az adatok (hirdetések,
tagságok, térképi objektumok) több nézetben, több állapotban jelennek meg. A kódbázis kizárólag
függvénykomponenseket és hookokat használ: a lokális állapotot a \texttt{useState}, a
mellékhatásokat a \texttt{useEffect}, a származtatott értékeket a \texttt{useMemo} kezeli, az
alkalmazásszintű hitelesítési állapotot pedig egy saját kontextus és a hozzá tartozó
\texttt{useAuth} hook szolgáltatja (\ref{chap:felhasznalok}.~fejezet).

A teljes kliensoldali kódbázis TypeScript nyelven készült~\cite{tshandbook}, szigorú
(\texttt{strict}) fordítóbeállításokkal, bekapcsolt holtkód-ellenőrzésekkel
(\texttt{noUnusedLocals}, \texttt{noUnusedParameters}) -- ezt a \texttt{tsconfig.app.json}
rögzíti. A statikus típusosság a projektben elsősorban a szerver--kliens határfelület
dokumentálását szolgálja: a backend válaszformátumait a \texttt{src/types} modul interfészei
tükrözik, így az API-t érintő eltérések már fordítási időben kiderülnek. Mivel a két oldal
típusai között nincs generált kapcsolat, az összhang fenntartása fegyelmet igényel; ez a
választott egyszerűség ára.

\subsection{Vite}

A fejlesztői kiszolgálót és a produkciós csomagolást a Vite adja~\cite{vitedocs}. Fejlesztés
közben a Vite natív ES-modulokat szolgál ki, így az indulás és a módosítás utáni frissítés
(hot module replacement) másodperc alatti; a produkciós build Rollup-alapú. A konfiguráció
(\texttt{vite.config.ts}) szándékosan minimális: a React-bővítmény és a Tailwind-bővítmény
mellett egyetlen útvonal-aliast definiál (\texttt{@} $\rightarrow$ \texttt{./src}), amely a
mély relatív importútvonalakat váltja ki, és amelyet a tesztkonfiguráció is örököl
(\ref{chap:teszteles}.~fejezet). A Vite \texttt{import.meta.env} mechanizmusa kezeli a
környezeti változókat is -- így jut el a kliensbe az API-alapcím és a külső szolgáltatások két
kulcsa (\ref{sec:teszt-uzemeltetes}.~alfejezet).

\subsection{TanStack Query (React Query)}

A szerverállapot kezelésére a rendszer a TanStack Query könyvtárat
használja~\cite{tanstackquery}. A könyvtár a szerverről származó adatokat deklaratívan,
kulcsokkal azonosított lekérdezésekként kezeli: gondoskodik a gyorsítótárazásról, az
állapotjelzésekről (betöltés, hiba) és az érvénytelenítés utáni újratöltésről. Ez a projektben
azt jelenti, hogy az oldalkomponensek nem írnak kézi adatbetöltő logikát: a \texttt{useQuery}
hook a megfelelő API-modult hívja, a mutációk pedig sikeres írás után az érintett kulcsokat
érvénytelenítik, mire a felület automatikusan frissül. A globális alapértelmezéseket -- egyperces
frissességi ablak, kikapcsolt fókusz-újratöltés és újrapróbálkozás -- a megosztott
\texttt{queryClient} példány rögzíti (\ref{sec:alapok-frontend}.~alfejezet). A kliensoldali
(nem szerverről származó) állapot ezzel szemben szándékosan egyszerű eszközökkel,
komponensállapottal és kontextussal van kezelve; külön állapotkezelő könyvtárra (Redux,
Zustand) a projekt méreténél fogva nem volt szükség.

\subsection{React Router, axios és a megjelenési réteg}

A kliensoldali útvonalkezelést a React Router 6-os verziója végzi~\cite{reactrouter}
(\texttt{createBrowserRouter}, beágyazott elrendezés-útvonalak, dinamikus szegmensek -- a teljes
útvonaltáblát a~\ref{sec:alapok-frontend}.~alfejezet közli). A HTTP-kommunikáció az axios
könyvtárra épül~\cite{axiosdocs}, amelynek elfogó- (interceptor-) mechanizmusa teszi lehetővé a
hitelesítés és a hibakezelés egyetlen pontra koncentrálását.

A megjelenés a Tailwind CSS 4-es verziójára épül~\cite{tailwinddocs}: a stílusok
segédosztályokként, közvetlenül a JSX-ben íródnak, a témát (színek, betűtípusok) pedig a
4-es verzió új, CSS-alapú \texttt{@theme} direktívája definiálja
(\ref{sec:alapok-frontend}.~alfejezet). A segédosztály-alapú megközelítést kiegészíti néhány
kis könyvtár: a \texttt{class-variance-authority} és a \texttt{tailwind-merge} a
variánsalapú UI-primitívek (gombok) megvalósítását segíti, a Radix UI \texttt{Slot}
komponense a gombok polimorf viselkedését adja, az ikonokat a \texttt{lucide-react} készlet
szolgáltatja, a chatbot válaszainak formázott megjelenítését pedig a \texttt{react-markdown}
végzi.

\section{A backend technológiái}
\label{sec:tech-backend}

\subsection{ASP.NET Core 8 Web API}

A szerveroldal az ASP.NET Core 8 keretrendszerre épül~\cite{aspnetcoredocs}, kizárólag Web
API-ként (nézetgenerálás nélkül). A keretrendszer három tulajdonsága bizonyult meghatározónak a
projektben. Az első a beépített \emph{függőséginjektálás}: a szolgáltatások életciklusát a
konténer kezeli, a kontrollerek és szolgáltatások konstruktoron át kapják függőségeiket, ami a
tesztelhetőség alapja (\ref{chap:teszteles}.~fejezet). A második a \emph{middleware-csővezeték},
amelyben a keresztvágó feladatok (CORS, hitelesítés, jogosultságkezelés, statikus fájlok)
deklaráltan, jól meghatározott sorrendben futnak (\ref{sec:alapok-backend}.~alfejezet). A
harmadik az \texttt{[ApiController]} konvenciókészlet: automatikus modellvalidáció,
egységes hibaválaszok és attribútumalapú útvonalillesztés. A hitelesítési igényeket az
ASP.NET Core Identity és a JWT-bearer middleware fedi le~\cite{identitydocs, rfc7519};
a fejlesztést a Swashbuckle által generált OpenAPI-leírás és Swagger-felület támogatja.

A megvalósítás a C\# nyelv modern eszközeire támaszkodik: a teljes adatelérés és minden
kontrollerakció aszinkron (\texttt{async}/\texttt{await}), ami a szálkészlet hatékony
kihasználását adja I/O-intenzív terhelés alatt; a bekapcsolt \emph{nullable referenciatípusok}
fordítási időben kényszerítik ki a hiányzó értékek kezelését; a rekordtípusok a kis,
változtathatatlan adathordozók (seed-definíciók, DTO-k) tömör leírását teszik lehetővé.

\subsection{Entity Framework Core}

Az adatelérést az Entity Framework Core 8 objektum-relációs leképező
végzi~\cite{efcoredocs}, Code First megközelítésben: a séma forrása a C\# entitásmodell, az
evolúciót migrációk írják le, a lekérdezések pedig LINQ-kifejezésekként készülnek, amelyeket a
szolgáltató SQL-re fordít. A projekt az EF Core képességeinek széles sávját használja: a
Fluent API-s modellkonfigurációt (összetett kulcsok, alternatív kulcsok, törlési viselkedések,
ellenőrző megszorítások -- \ref{sec:alapok-adat}.~alfejezet), a változáskövetést az írási
műveleteknél, az \texttt{AsNoTracking} csak olvasó lekérdezéseket a listázásoknál, az
\texttt{Include}-alapú betöltést a kapcsolt adatokhoz, valamint a kiszolgálóoldalra forduló
lapozást és szűrést (\ref{chap:hirdetesek}.~fejezet). Az ORM választása a kézi SQL-lel szemben
a típusbiztonságot és a sémával együtt verziózott migrációkat hozta; az absztrakció árát --
a generált lekérdezések feletti részleges kontrollt -- a projekt adatmennyiségei mellett nem
ellensúlyozza kézi optimalizálási igény.

\section{Microsoft SQL Server}
\label{sec:tech-mssql}

Az adatok a Microsoft SQL Server relációs adatbázis-kezelőben
tárolódnak~\cite{sqlserverdocs}. A választást a .NET-ökoszisztémával való szoros integráció
(első osztályú EF Core-szolgáltató) és a fejlesztést egyszerűsítő LocalDB-kiadás indokolta: a
LocalDB igény szerint induló, telepítést alig igénylő folyamatként futtatja a teljes értékű
SQL Server-motort, így a fejlesztői gépen nincs szükség adatbázis-kiszolgáló üzemeltetésére. A
\emph{VadászNet} nevű adatbázist a kapcsolati karakterlánc jelöli ki; a sémát és a kezdeti
tartalmat a backend induláskor, migrációkkal és adatfeltöltéssel maga állítja elő
(\ref{sec:alapok-adat}.~alfejezet), így az adatbázis a verziókezelt kódból bármikor
reprodukálható. A rendszer kihasznál néhány kiszolgálóoldali képességet is: az időbélyegek
\texttt{sysdatetime()} alapértékét, az egyedi indexeket és az ellenőrző megszorításokat.

\section{Külső szolgáltatások}
\label{sec:tech-kulso}

A rendszer három külső szolgáltatást integrál, kétféle mintával. A \emph{Google Maps JavaScript
API}~\cite{gmapsdocs} a térképes funkciók alapja: a kliens a \texttt{@react-google-maps/api}
csomag React-komponensein keresztül jeleníti meg a térképet, a jelölőket és a rajzolóeszközt;
a betöltés kliensoldali API-kulcshoz kötött (\ref{chap:terkep}.~fejezet). Az \emph{OpenAI Chat
Completions API}~\cite{openaidocs} a vadászati asszisztens nyelvi modelljét (gpt-4o)
szolgáltatja; ezt kizárólag a backend hívja, szerveroldali kulccsal, saját végpont mögött
(\ref{chap:chatbot}.~fejezet). Az \emph{OpenWeather API}~\cite{openweatherdocs} aktuális
időjárási adatokat ad településnév alapján; ezt -- a másik két szolgáltatástól eltérően -- a
kliens közvetlenül hívja, ami egyszerűbb megvalósítás, de a kulcs kitettségével jár
(\ref{chap:chatbot}.~fejezet). A kétféle integrációs minta (szerveroldali proxy, illetve
közvetlen klienshívás) összevetése a~\ref{sec:arch-attekintes}.~alfejezetben, architekturális
szinten is megjelenik.

% !TeX root = ../main.tex
% 3. fejezet — Rendszerarchitektúra
% A ftsrg/thesis-template-latex sablonba \include-dal illesztendő tartalomfájl.
% Hivatkozott források: Backend/HUNting.API/Program.cs, Controllers/*, Services/*,
%   Frontend/HUNting/src/lib/apiClient.ts, services/weatherService.ts, pages/Map/MapPage.tsx,
%   components/Chatbot.tsx, thesis/planning/feature_matrix.md

\chapter{Rendszerarchitektúra}
\label{chap:architektura}

Ez a fejezet madártávlatból mutatja be a HUNting felépítését: a rétegeket és felelősségeiket, a
rétegek közötti kommunikáció szabályait, valamint azokat a tervezési döntéseket, amelyek a
rendszer jelenlegi formáját meghatározzák. A cél a mérnöki szempontból pontos minősítés: a
HUNting \emph{rétegzett monolit} -- nem mikroszolgáltatás-architektúra és nem is a Clean
Architecture értelmében vett, szigorúan elválasztott rétegrendszer --, és a fejezet azt is
megindokolja, hogy a projekt méretéhez és céljaihoz miért ez az ésszerű választás. A rétegek
\emph{belső} felépítését (adatbázis-kontextus, DI, middleware, kliensváz) a
\ref{chap:alapok}.~fejezet részletezi; itt a rétegek \emph{közötti} viszonyokra koncentrálunk.

\section{A háromrétegű architektúra áttekintése}
\label{sec:arch-attekintes}

A rendszer klasszikus háromrétegű felépítést követ, amelyet
a~\ref{fig:architektura}.~ábra szemléltet. A \emph{megjelenítési réteg} egy React-alapú
egyoldalas alkalmazás (single-page application, SPA): a böngésző egyszer tölti le a statikus
HTML-, JavaScript- és CSS-állományokat, ezt követően minden képernyőváltás kliensoldali
útvonalkezeléssel, minden adatművelet aszinkron HTTP-hívással történik. Az \emph{üzleti réteg}
egy ASP.NET Core 8 Web API, amely állapotmentes REST-végpontokon szolgálja ki a klienst: itt
történik a hitelesítés és a jogosultság-ellenőrzés, az üzleti szabályok érvényesítése és a
feltöltött fájlok kezelése. Az \emph{adatréteg} egy Microsoft SQL Server adatbázis (fejlesztői
környezetben LocalDB), amelyet a backend kizárólag az Entity Framework Core objektum-relációs
leképezőn keresztül ér el -- kézzel írt SQL a kódbázisban nem fordul elő.

% TODO: A-01 ábra (figures/a-01-architektura.pdf) — React SPA ↔ REST/JSON ↔ ASP.NET Core API ↔
% EF Core ↔ MSSQL; külső API-k: OpenAI (szerveroldalról), Google Maps + OpenWeather (kliensről).
\begin{figure}[htbp]
	\centering
	\fbox{\parbox[c][7cm][c]{0.9\textwidth}{\centering A-01: A HUNting háromrétegű
	architektúrája a külső szolgáltatásokkal (elkészítendő)}}
	\caption{A rendszer architektúrája: rétegek és külső integrációk}
	\label{fig:architektura}
\end{figure}

A repository ezt a felépítést két fő projektben tükrözi: a \texttt{Backend/HUNting.API} az
API-t és az adatelérést, a \texttt{Frontend/HUNting} a kliensalkalmazást tartalmazza; mindkettő
mellett saját tesztprojekt, illetve tesztkönyvtár áll (\ref{chap:teszteles}.~fejezet). A két
projekt között fordítási idejű függőség nincs: a kapcsolatot kizárólag a HTTP-határfelület és
a JSON-formátum jelenti, a kliensoldali TypeScript-típusok a backend válaszformátumait kézzel
tükrözik.

A képet a külső szolgáltatások teszik teljessé, és itt egy architekturális szempontból fontos
aszimmetria figyelhető meg. Az OpenAI Chat Completions API-t a backend hívja (a chatbot-kérések
a saját \texttt{/api/Chat} végponton keresztül, szerveroldali API-kulccsal jutnak el az OpenAI
felé), a Google Maps JavaScript API-t és az OpenWeather API-t viszont közvetlenül a kliens éri
el, kliensoldali kulcsokkal. Ez utóbbi egyszerűbb megvalósítást adott, ára azonban a kulcsok
böngészőbeli kitettsége; a kompromisszumot a~\ref{chap:chatbot}.~fejezet elemzi.

A rendszer szereplőit és fő használati eseteit a~\ref{fig:usecase}.~ábra foglalja össze. A
\emph{vendég} böngészheti a hirdetéseket, a társaságkatalógust és az időjárásoldalt, valamint
használhatja a chatbotot; a \emph{vadász} (Person fiók) emellett hirdetést ad fel, társasági
tagságot igényel, naplót vezet és térképi objektumokat kezel; a \emph{társasági fiók} (Society)
a tagfelvételt bírálja el, bejegyzéseket és tagokat kezel; az \emph{adminisztrátor} a teljes
felhasználó- és hirdetésállomány felett rendelkezik.

% TODO: U-01 ábra (figures/u-01-usecase.pdf) — Use Case diagram: Vendég, Vadász (Person),
% Társasági fiók (Society), Adminisztrátor aktorok a 6 funkcióterület fő eseteivel.
\begin{figure}[htbp]
	\centering
	\fbox{\parbox[c][7cm][c]{0.9\textwidth}{\centering U-01: A HUNting használati esetei
	(elkészítendő)}}
	\caption{A rendszer szereplői és fő használati esetei}
	\label{fig:usecase}
\end{figure}

\section{REST API és adatfolyamok}
\label{sec:arch-rest}

A kliens és a szerver közötti határfelület a REST elveit követi~\cite{fielding2000}. Az
erőforrásokat főnévi útvonalak azonosítják az egységes \texttt{api/[controller]} sémában
(\texttt{/api/Ads}, \texttt{/api/Societies}, \texttt{/api/HuntingLogs}~stb.), a műveletet pedig
a HTTP-ige hordozza: a \texttt{GET} lekérdez, a \texttt{POST} létrehoz (illetve alárendelt
akciókat indít, például \texttt{POST /api/Societies/\{id\}/join}), a \texttt{PUT} teljes
módosítást végez, a \texttt{DELETE} töröl. A válaszok a szabványos státuszkódokkal élnek:
sikeres lekérdezésre 200, létrehozásra 201 (\texttt{Created}, a létrejött erőforrásra mutató
hivatkozással), tartalom nélküli sikerre 204; a hibás kérés 400-at, a hiányzó vagy érvénytelen
token 401-et, az elégtelen jogosultság 403-at, a nem létező erőforrás 404-et kap. A
reprezentáció mindkét irányban JSON; a szerializáció globális beállításait (karakterlánc-alapú
enumok, körkörös hivatkozások elvágása, \texttt{null} mezők elhagyása) a
\ref{chap:alapok}.~fejezet mutatta be.

Egy tipikus kérés életútját a~\ref{fig:request-flow}.~ábra követi végig. A kliensoldali
funkciómodul az \texttt{apiClient} példányon indítja a hívást, amelynek kérés-elfogója befűzi
a bearer-tokent. A kérés a szerveren a middleware-csővezetéken halad át: a CORS-ellenőrzés után
a hitelesítési middleware validálja a tokent és felépíti a kérés identitását, a
jogosultságkezelés kiértékeli az \texttt{[Authorize]} szabályokat, majd az útvonalillesztés a
megfelelő kontrollerakciót hívja. Az akció az erőforrás-szintű ellenőrzések után a
szolgáltatásrétegen keresztül az EF Core-ral fordul az adatbázishoz, az eredmény pedig
JSON-válaszként utazik vissza, ahol a React Query gyorsítótárba kerül, és a komponensek
deklaratívan megjelenítik. Hiba esetén a kliens válasz-elfogója lép közbe: a 401-es választ
globális kijelentkeztetésként kezeli, minden mást egységes hibaüzenetté normalizál.

% TODO: A-02 ábra (figures/a-02-keres-eletut.pdf) — folyamatábra: apiClient interceptor → CORS →
% authentikáció → autorizáció → controller → service → EF Core → MSSQL → JSON válasz → React Query.
\begin{figure}[htbp]
	\centering
	\fbox{\parbox[c][6.5cm][c]{0.9\textwidth}{\centering A-02: Egy REST-kérés életútja a
	kliens-elfogótól az adatbázisig és vissza (elkészítendő)}}
	\caption{Egy hitelesített REST-kérés feldolgozási folyamata}
	\label{fig:request-flow}
\end{figure}

Az adatfolyamok szempontjából a rendszer három forgalomtípust bonyolít. A \emph{strukturált
adatforgalom} a fenti JSON-séma szerint zajlik. A \emph{bináris forgalom} a képfeltöltésekre és
-letöltésekre korlátozódik: a feltöltés \texttt{multipart/form-data} formátumú POST-kérés, a
letöltés pedig a statikus fájl-middleware-en keresztül, a JSON-válaszokban kapott URL-ek alapján
történik. Végül a \emph{külső forgalom} a chatbot-kérések továbbítása az OpenAI felé
(szerveroldalon), illetve a térkép- és időjárásadatok lekérése (kliensoldalon).

\section{Tervezési döntések és kompromisszumok}
\label{sec:arch-dontesek}

\paragraph{Monolit a mikroszolgáltatások helyett.} A backend egyetlen telepítési egység, közös
adatbázissal. Egy hat funkcióterületből álló, egy fejlesztő által készített alkalmazásnál a
mikroszolgáltatásokra bontás (szolgáltatásonkénti adatbázis, szolgáltatásközi kommunikáció,
elosztott tranzakciók) aránytalan üzemeltetési és fejlesztési többletköltséget jelentene,
miközben az előnyei -- független skálázás és telepítés -- itt nem jelentkeznének. A monoliton
belül a funkcióterületek a kontroller--szolgáltatás--modell hármasok mentén így is jól
elkülönülnek, ami a későbbi kiemelést nem zárja ki.

\paragraph{Rétegzett monolit, nem Clean Architecture.} A rétegek elválasztása pragmatikus, nem
dogmatikus: az entitásosztályok egyszerre szolgálnak adatbázis-leképezésként és -- több végpontnál
-- API-válaszformátumként, a szolgáltatásréteg pedig nem mindenhol van jelen (a hírek kontrollere
közvetlenül a kontextust használja). Külön tartományi réteg, interfészprojekt vagy
DTO-leképező keretrendszer (például AutoMapper) nincs: ahol a válasz formátuma eltér az
entitástól (például a hitelesítési válaszok), ott kézi leképezés történik. Ez a döntés
csökkenti az absztrakciós rétegek számát és a kód mennyiségét, cserébe a séma és az API
összefonódik: egy entitásmódosítás közvetlenül kihat a válaszformátumra. A dolgozat ezt az
őszinte minősítést következetesen tartja: ahol egy fejezet a rétegzéstől való eltérést talál,
ott azt jelzi (így a részleges repository-mintát a~\ref{sec:alapok-backend}.~alfejezet
tárgyalja~\cite{fowler2002}).

\paragraph{Oldalszintű védelem a kliensen.} A kliensoldali hozzáférés-védelem nem központi
útvonalőrrel, hanem az egyes oldalkomponensekben valósul meg
(\ref{sec:alapok-frontend}.~alfejezet). Ez a megközelítés egyszerű és oldalanként testre
szabható (eltérő átirányítási célok, szerepkörfüggő nézetek), hátránya a szabályok szétszórtsága.
Hangsúlyozandó ugyanakkor az architektúra alapelve: a kliensoldali védelem kizárólag felhasználói
élményt szolgál, a tényleges hozzáférés-szabályozás minden esetben a szerveren történik -- a
kliens megkerülésével indított kérések is a JWT-validáción és a szerepkör-, illetve
tulajdonos-ellenőrzéseken buknak el.

\paragraph{A funkcióközpontú dokumentálás indoklása.} A dolgozat szerkezete maga is tervezési
döntés. Mivel a rendszer értéke a funkcióiban rejlik, és minden funkció ugyanazon a három
rétegen ível át, a dokumentáció nem rétegenként (külön adatbázis-, backend- és
frontend-fejezetekben), hanem funkciónként halad: az 5--10.~fejezet mindegyike az adatmodelltől
a felhasználói felületig mutat be egy-egy funkciót. Így az összetartozó döntések egy helyen
olvashatók, az ismétlődő infrastruktúra pedig a \ref{chap:alapok}.~fejezetben, egyszer szerepel.

% !TeX root = ../main.tex
% 4. fejezet — Közös rendszeralapok
% A ftsrg/thesis-template-latex sablonba \include-dal illesztendő tartalomfájl.
% Hivatkozott források: Backend/HUNting.API/Program.cs, Data/HuntingContext.cs, Data/SeedData.cs,
%   Configuration/JwtOptions.cs, Services/UploadStorage.cs, Controllers/ApiControllerBase.cs,
%   Services/AdService.cs, SocietyService.cs, Controllers/* (végponttábla, feltöltési minta),
%   Migrations/*, appsettings.json, Properties/launchSettings.json,
%   Frontend/HUNting/src/main.tsx, App.tsx, routes/appRouter.tsx, layouts/AppLayout.tsx,
%   AuthLayout.tsx, lib/apiClient.ts, queryClient.ts, authStorage.ts, context/AuthContext.tsx,
%   components/layout/AppHeader.tsx, index.css

\chapter{Közös rendszeralapok}
\label{chap:alapok}

A HUNting funkciói -- a felhasználókezeléstől a térképes modulig -- közös infrastruktúrára
épülnek: ugyanazt az adatbázis-kontextust, ugyanazt a hitelesítési middleware-t és ugyanazt a
kliensoldali HTTP-réteget használják. Ez a fejezet ezt a funkciófüggetlen alapréteget mutatja be
egyszer, összefüggően, hogy a dolgozat gerincét adó
\ref{chap:felhasznalok}--\ref{chap:chatbot}.~funkciófejezetek ismétlés nélkül, kereszthivatkozással
építhessenek rá. A tárgyalás a rendszer rétegeit alulról felfelé követi: először az adattárolás
alapjait (adatbázis-kontextus, migrációk, kezdeti adatfeltöltés), majd a backend közös elemeit
(alkalmazás-összeállítás, rétegzés, JWT-infrastruktúra, visszatérő megvalósítási minták), végül a
frontend vázát (provider-fa, útvonalkezelés, API-kommunikáció, megjelenés) ismerteti.

\section{Adattárolási alapok}
\label{sec:alapok-adat}

\subsection{A \texttt{HuntingContext} és az adatmodell}

Az adatelérés az Entity Framework Core \emph{Code First} megközelítésére épül: az adatbázisséma
egyetlen hiteles forrása a C\# entitásmodell, amelyből a séma migrációkkal áll elő. A fejlesztői
környezetben az adatbázis egy \emph{VadászNet} nevű SQL Server LocalDB-példány, amelyet a
konfiguráció \texttt{DefaultConnection} kapcsolati karakterlánca jelöl ki.

A rendszer egyetlen adatbázis-kontextusa a \texttt{HuntingContext}
(\texttt{Data/HuntingContext.cs}), amely nem a sima \texttt{DbContext}, hanem az
\texttt{IdentityDbContext<ApplicationUser, IdentityRole<int>, int>} osztályból származik: így az
ASP.NET Core Identity felhasználó- és szereptáblái ugyanabban a sémában élnek, mint az üzleti
entitások, a teljes modell pedig egységesen egész számú elsődleges kulcsokat használ. A kontextus
húsz \texttt{DbSet} tulajdonságot tesz közzé, amelyek tartományok szerint csoportosíthatók:

\begin{itemize}[nosep]
	\item \emph{Felhasználók:} \texttt{Users} (az Identity alaptábláját fedi le);
	\item \emph{Társaságok:} \texttt{Societies}, \texttt{SocietyRoles}, \texttt{SocietyMembers},
	      \texttt{SocietyMembershipRequests}, \texttt{SocietyPosts};
	\item \emph{Hirdetések:} \texttt{Ads}, \texttt{AdCategories}, \texttt{AdImages},
	      \texttt{AdStatuses};
	\item \emph{Hírek:} \texttt{NewsFeed}, \texttt{NewsVisibility};
	\item \emph{Vadásznapló:} \texttt{HuntingLogs};
	\item \emph{Térkép:} \texttt{MapVisibilities}, \texttt{MapIconTypes}, \texttt{MapMarkers},
	      \texttt{MapPolygons}, \texttt{MapPolygonPoints}, \texttt{MapPaths},
	      \texttt{MapPathPoints}.
\end{itemize}

Ezekhez az Identity további hat infrastruktúra-táblája (szerepek, szerep-hozzárendelések,
claimek, külső bejelentkezések, tokenek) társul, így az adatbázis összesen 26 táblából áll. Az
Identity alapértelmezett \texttt{AspNet*} előtagú táblaneveit a kontextus
\texttt{OnModelCreating} metódusa beszédesebb nevekre (\texttt{Users}, \texttt{Roles},
\texttt{UserRoles}, \texttt{UserClaims}, \texttt{UserLogins}, \texttt{RoleClaims},
\texttt{UserTokens}) képezi le.

A modellkonfiguráció a Fluent API-ra és az entitásokon elhelyezett DataAnnotations
attribútumokra épül; a visszatérő megoldások a következők.
A státusz- és láthatóságjellegű \emph{kódtáblák} (\texttt{AdStatus}, \texttt{NewsVisibility},
\texttt{MapVisibility}) karakterlánc-kódokkal (például \texttt{ACTIVE}, \texttt{PUBLIC},
\texttt{MEMBERS}) kapcsolódnak a hivatkozó entitásokhoz: az idegen kulcs nem a numerikus
azonosítóra, hanem a kódoszlopra mutat (a \texttt{HasPrincipalKey} beállítással), így az
adatbázis-rekordok és az API-válaszok önmagukban is olvashatók. A kapcsolótáblák
(\texttt{SocietyMember}, \texttt{SocietyMembershipRequest}) \emph{összetett elsődleges kulcsot}
kapnak a társaság- és felhasználó-azonosító párosából, ami sémaszinten zárja ki a duplikált
tagságot és kérelmet. Az üzleti egyediséget \emph{egyedi indexek} garantálják (felhasználói és
társasági e-mail-cím, társasági szerepnév, hirdetéskategória-név, valamint a térképi pontoknál a
poligon-, illetve útvonal-azonosító és a \texttt{Seq} sorszám párosa). A létrehozási időbélyegek
(\texttt{CreatedAt}, \texttt{JoinedAt}) adatbázisoldali \texttt{sysdatetime()} alapértéket kapnak.
Végül a hirdetéseknél egy \emph{ellenőrző megszorítás} (\texttt{CK\_Ads\_SingleOwner}) mondja ki,
hogy minden hirdetésnek pontosan egy tulajdonosa van: vagy felhasználó, vagy társaság, de sosem
mindkettő vagy egyik sem. A teljes adatmodellt a~\ref{fig:er-diagram}.~ábra foglalja össze; az
egyes tartományok entitásait a megfelelő funkciófejezetek adatmodell-alfejezetei részletezik.

% TODO: E-01 ábra (figures/e-01-er-diagram.pdf) — teljes ER-diagram, Mermaid erDiagram kódból
% generálható a HuntingContext.cs alapján; egész oldalas, fekvő tájolás javasolt.
\begin{figure}[htbp]
	\centering
	\fbox{\parbox[c][9cm][c]{0.95\textwidth}{\centering E-01: A VadászNet adatbázis ER-diagramja
	(elkészítendő) -- Identity-mag + társasági, hirdetési, hír-, napló- és térképi tartomány}}
	\caption{A VadászNet adatbázis entitás--kapcsolat diagramja}
	\label{fig:er-diagram}
\end{figure}

\subsection{Törlési viselkedések}

A kapcsolatok törlési viselkedését a kontextus kapcsolatonként, kifejezetten rögzíti; ez a
relációs integritás szempontjából a modell egyik legfontosabb tervezési döntéssorozata.
A~\ref{tab:delete-behavior}.~táblázat a három alkalmazott viselkedés szerint csoportosítva
tekinti át a kapcsolatokat.

\begin{table}[htbp]
	\centering
	\caption{Törlési viselkedések a \texttt{HuntingContext} kapcsolataiban}
	\label{tab:delete-behavior}
	\begin{tabular}{|l|p{10.2cm}|}
		\hline
		\textbf{Viselkedés} & \textbf{Kapcsolatok} \\
		\hline
		\texttt{Cascade} &
		tagság és kérelem a társaság felől; tagság a felhasználó felől; társasági bejegyzés a
		társaság és a szerző felől; hirdetés a tulajdonos felhasználó, illetve társaság felől;
		hirdetéskép a hirdetés felől; hír a társaság és a szerző felől; naplóbejegyzés a
		felhasználó felől; térképi objektum a társaság felől; poligon- és útvonalpont a
		poligon, illetve útvonal felől \\
		\hline
		\texttt{Restrict} &
		tagság a társasági szerep felől; kérelem a kérelmező felhasználó felől; hirdetés a
		kategória és a státusz felől; hír a láthatósági kód felől; térképi objektum a
		létrehozó felhasználó, az ikontípus és a láthatósági kód felől \\
		\hline
		\texttt{SetNull} &
		a társaság \texttt{OwnerUserId} kulcsa a tulajdonos felhasználó törlésekor; a kérelem
		\texttt{ProcessedByUserId} kulcsa az elbíráló törlésekor \\
		\hline
	\end{tabular}
\end{table}

A mintázat következetes logikát tükröz. A \emph{tartalmazási} viszonyokban (a hirdetés képei, a
poligon pontjai, a társaság bejegyzései) a szülő törlése kaszkádban viszi magával a gyermekeket,
mert azok a szülő nélkül értelmezhetetlenek. A \emph{hivatkozási} viszonyokban -- ahol a
hivatkozott rekord megosztott kódtábla- vagy törzsadat-elem -- a \texttt{Restrict} akadályozza meg,
hogy egy még használatban lévő kategória, státusz vagy szerep eltűnjön a hivatkozó rekordok alól.
A \texttt{SetNull} ott jelenik meg, ahol a kapcsolat elvesztése nem érvényteleníti a rekordot:
a társaság a tulajdonosa nélkül is létező szervezet marad, az elbírált kérelem pedig az elbíráló
fiókjának megszűnése után is őrzi az állapotát. Megjegyzendő ugyanakkor egy aszimmetria: a térképi
objektumoknál a létrehozó felhasználóra mutató kulcs \texttt{Restrict} viselkedésű, vagyis a
felhasználó mindaddig nem törölhető, amíg térképi objektuma van -- miközben naplóbejegyzései és
hirdetései kaszkádban törlődnének. Ez a gyakorlatban azt jelenti, hogy az adminisztrátori
felhasználótörlés térképi tartalommal rendelkező fióknál adatbázishibára fut; a \texttt{Restrict}
viselkedésre a térképfejezet adatmodell-alfejezete (\ref{sec:map-adatmodell}.~alfejezet) is kitér.

\subsection{Migrációk}

A séma evolúcióját öt EF Core migráció írja le, amelyek egyben a fejlesztés idővonalát is
dokumentálják:

\begin{enumerate}[nosep]
	\item \texttt{InitialIdentity} (2025.~október) -- az Identity-alapú felhasználókezelés,
	      valamint a társasági, hirdetési, hír- és naplótartomány kezdeti sémája;
	\item \texttt{AddMapEntities} (2025.~november) -- a térképi entitások (láthatóságok,
	      ikontípusok, jelölők, poligonok, útvonalak és pontjaik) bevezetése;
	\item \texttt{SocietyAccountsAndPosts} (2025.~november) -- a társasági fiókok és a
	      társasági bejegyzések támogatása;
	\item \texttt{RenameIdentityTables} (2025.~december) -- az \texttt{AspNet*} táblanevek
	      lecserélése a beszédes nevekre;
	\item \texttt{AddDescriptionToMapPolygon} (2026.~március) -- leírásmező a térképi
	      poligonokhoz.
\end{enumerate}

A fejlesztői munkafolyamatban a migrációk a \texttt{dotnet ef migrations add} paranccsal
készülnek, alkalmazásukhoz azonban nincs szükség kézi \texttt{dotnet ef database update}
futtatásra: az alkalmazás indulásakor a kezdeti adatfeltöltés első lépése a
\texttt{Database.MigrateAsync()} hívás, amely a még hiányzó migrációkat automatikusan érvényesíti.
Ez a megoldás fejlesztői és demonstrációs környezetben kényelmes -- az adatbázis a backend első
indításával üzemkész --, éles, többpéldányos környezetben ugyanakkor a telepítéshez kötött,
ellenőrzött migrációfuttatás lenne indokolt.

\subsection{Kezdeti adatfeltöltés}

A \texttt{SeedData.SeedAsync} statikus metódus (\texttt{Data/SeedData.cs}) az alkalmazás minden
indulásakor lefut, és a migrációk érvényesítése után feltölti az adatbázist a működéshez és a
bemutatáshoz szükséges adatokkal. A feltöltés kétféle adatot érint. A \emph{törzsadatok} a rendszer
működéséhez szükségesek: a három alkalmazásszintű szerep (\emph{Administrator},
\emph{SocietyManager}, \emph{Hunter}), négy hirdetésstátusz, három-három hír- és térképláthatósági
kód, négy térképi ikontípus, négy társasági szerep (Vezető, Gazda, Vadőr, Tag) és tizenkét
hirdetéskategória. A \emph{demonstrációs adatok} a felületek bemutatását szolgálják: tizenegy
felhasználó (egy adminisztrátor és tíz vadász), nyolc vadásztársaság, tagsági kapcsolatok, hét
hirdetés képekkel, két hírbejegyzés és négy vadásznapló-bejegyzés.

A feltöltés kulcstulajdonsága az \emph{idempotencia}: minden rekord beszúrását egy természetes
kulcsra (e-mail-címre, névre, kódra vagy címre) szűrő létezés-ellenőrzés előzi meg, így az
ismételt indítás nem duplikál adatot, a hiányzó elemeket viszont pótolja.
A~\ref{lst:seed}.~kódrészlet ezt a mintát szemlélteti a társaságok feltöltésén.

\begin{lstlisting}[language={[Sharp]C}, caption={Idempotens adatfeltöltés a \texttt{SeedData} osztályban (kivonat)}, label={lst:seed}]
private static async Task EnsureSocietiesAsync(HuntingContext context)
{
    var societies = new[]
    {
        new Society { Name = "Bakonyi Vadasztarsasag",
                      Email = "kapcsolat@bakony-vt.example.hu", ... },
        // ... tovabbi het tarsasag ...
    };

    foreach (var society in societies)
    {
        if (!await context.Societies.AnyAsync(s => s.Email == society.Email))
        {
            context.Societies.Add(society);
        }
    }
}
\end{lstlisting}

A megoldás korlátaival együtt teljes a kép. A kezdő jelszavak (köztük az adminisztrátoré) a
forráskódban szerepelnek, ami kizárólag fejlesztői környezetben elfogadható. A
\emph{SocietyManager} szerep létrejön ugyan, de a feltöltött felhasználók egyike sem kapja meg: a
nyolc demonstrációs társaság közvetlenül, tulajdonos felhasználói fiók nélkül kerül az adatbázisba
(az \texttt{OwnerUserId} üresen marad, a kötelező \texttt{PasswordHash} mezőbe helykitöltő érték
kerül), társasági fiók így csak a regisztrációs folyamaton keresztül keletkezik
(\ref{chap:felhasznalok}.~fejezet). A hirdetésképek és a bejegyzésképek külső (Unsplash) URL-ekre
mutatnak, nem a lokális feltöltési könyvtárra.

\section{Backend-alapok}
\label{sec:alapok-backend}

\subsection{Az alkalmazás összeállítása}

A backend egyetlen belépési pontja a \texttt{Program.cs}, amely a .NET 8 minimál hoszting
modelljét követi: felső szintű utasításokkal, külön \texttt{Startup} osztály nélkül építi fel a
szolgáltatáskonténert és a HTTP-csővezetéket. A~\ref{lst:program-bootstrap}.~kódrészlet a fájl
szerkezetét mutatja a lényeges regisztrációkra és a middleware-sorrendre szűkítve.

\begin{lstlisting}[language={[Sharp]C}, caption={DI-regisztrációk és middleware-sorrend a \texttt{Program.cs}-ben (kivonat)}, label={lst:program-bootstrap}]
builder.Services.AddDbContext<HuntingContext>(opts =>
    opts.UseSqlServer(builder.Configuration
        .GetConnectionString("DefaultConnection")));

builder.Services.AddIdentityCore<ApplicationUser>(...)
    .AddRoles<IdentityRole<int>>()
    .AddEntityFrameworkStores<HuntingContext>()
    .AddSignInManager();

builder.Services.AddAuthentication(...).AddJwtBearer(...);
builder.Services.AddCors(...);

builder.Services.AddScoped<IAdService, AdService>();
builder.Services.AddScoped<ITokenService, JwtTokenService>();
// ... tovabbi service-regisztraciok ...
builder.Services.AddHttpClient<IOpenAiService, OpenAiService>();

var app = builder.Build();

if (app.Environment.IsDevelopment()) { app.UseSwagger(); app.UseSwaggerUI(); }
app.UseHttpsRedirection();
app.UseStaticFiles();      // wwwroot/uploads kiszolgalasa
app.UseCors("Frontend");
app.UseAuthentication();
app.UseAuthorization();
app.MapControllers();

await SeedData.SeedAsync(app.Services);
app.Run();
\end{lstlisting}

A szolgáltatásregisztrációk három csoportra oszthatók. Az \emph{infrastruktúra-szolgáltatások}
közé tartozik az adatbázis-kontextus, az Identity (jelszókezelés, szerepkezelés,
\texttt{SignInManager}) és a JWT-bearer hitelesítés. Az Identity jelszószabályzata szándékosan
enyhített: hat karakteres minimumhossz mellett sem kis- és nagybetűt, sem számjegyet, sem
speciális karaktert nem követel meg -- ez a demonstrációs célú használatot könnyíti, éles
üzemben szigorítandó. Az \emph{üzleti szolgáltatások} scoped életciklussal regisztrált
interfész--implementáció párok (\texttt{IAdService}, \texttt{IUserService},
\texttt{ISocietyService}, \texttt{IMapService}, \texttt{ITokenService}). A \emph{külső integrációt} az
\texttt{IOpenAiService} képviseli, amely típusos \texttt{HttpClient}-ként
(\texttt{AddHttpClient}) regisztrálódik, így a keretrendszer kezeli a HTTP-kapcsolatok
életciklusát (\ref{chap:chatbot}.~fejezet).

A kontrollerek JSON-szerializációja négy globális beállítást kap: a navigációs tulajdonságok
okozta körkörös hivatkozásokat a \texttt{ReferenceHandler.IgnoreCycles} vágja el, a \texttt{null}
értékű mezők kimaradnak a válaszokból, a szerializálási mélység legfeljebb 16 szint, a felsorolt
típusok pedig a \texttt{JsonStringEnumConverter} révén nevükkel (nem numerikus értékükkel)
utaznak -- így lesz a fióktípus a kliens számára \texttt{"Person"} vagy \texttt{"Society"}. A
\texttt{Frontend} nevű CORS-házirend a Vite fejlesztői kiszolgáló
(\texttt{http(s)://localhost:5173}) származási helyét engedélyezi tetszőleges fejléccel és
metódussal; a házirend a csővezetékben a hitelesítés \emph{előtt} fut le, így a böngészői
előkérések (preflight) tokenek nélkül is sikeresek. A middleware-sorrend további tagjai a
HTTPS-átirányítás, a statikus fájlok kiszolgálása (ez teszi elérhetővé a feltöltött képeket
tartalmazó \texttt{wwwroot/uploads} könyvtárat), majd a hitelesítés és a jogosultságkezelés. A
fejlesztői környezetben a Swagger (OpenAPI) felület is bekapcsolódik, amely a teljes
végpontkészletet dokumentálja és próbahívásokra is alkalmas
(\ref{fig:swagger}.~ábra).

% TODO: SC-17 képernyőkép (figures/sc-17-swagger.png) — Development módban, a végpontlistával.
\begin{figure}[htbp]
	\centering
	\fbox{\parbox[c][5.5cm][c]{0.85\textwidth}{\centering SC-17: A Swagger UI a HUNting.API
	végpontlistájával (elkészítendő)}}
	\caption{A generált Swagger-felület fejlesztői környezetben}
	\label{fig:swagger}
\end{figure}

\subsection{Rétegzés: kontrollerek és szolgáltatások}

A backend rétegzése a kontroller--szolgáltatás--adatelérés hármasra épül: a kontroller a
HTTP-szintű feladatokat (útvonalillesztés, státuszkódok, jogosultság-ellenőrzés) látja el, az
üzleti logika az interfész mögé rejtett szolgáltatásosztályokban él, az adatelérés pedig az EF
Core kontextuson keresztül történik. A~\ref{fig:komponens}.~ábra a kontrollerek, szolgáltatások
és a kontextus tényleges függőségi viszonyait mutatja.

% TODO: A-03 ábra (figures/a-03-komponensek.pdf) — komponensdiagram a Program.cs regisztrációi
% alapján: 8 controller → service-interfészek → HuntingContext; typed HttpClient → OpenAI.
\begin{figure}[htbp]
	\centering
	\fbox{\parbox[c][7cm][c]{0.9\textwidth}{\centering A-03: Backend-komponensdiagram
	(elkészítendő) -- kontrollerek, szolgáltatások, \texttt{HuntingContext}, külső OpenAI API}}
	\caption{A backend komponensei és függőségeik}
	\label{fig:komponens}
\end{figure}

Fontos azonban rögzíteni, hogy a rétegzés nem egységes. A szolgáltatásosztályok -- így az
\texttt{AdService} és a \texttt{SocietyService} -- közvetlenül a \texttt{HuntingContext}-tel
dolgoznak, külön repository-absztrakció nincs; a hírek és a vadásznapló kontrollere pedig
szolgáltatásréteg nélkül, közvetlenül a kontextusra épül. A kontrollerekben ismétlődő
segédlogika ugyanakkor két közös helyre szervezett: az \texttt{ApiControllerBase} ősosztály
adja a bejelentkezett felhasználó azonosítójának kiolvasását, valamint a szerep- és
fióktípus-ellenőrzéseket (\texttt{GetUserId}, \texttt{IsAdministrator},
\texttt{IsSocietyAccount}), a statikus \texttt{UploadStorage} segédosztály pedig a
képfeltöltések fájlrendszer-műveleteit végzi (lásd a közös megvalósítási mintáknál). Ahol egy
funkciófejezet a rétegzéstől eltér, ott azt külön jelzi.

\subsection{A JWT-infrastruktúra}

A hitelesítés állapotmentes, JSON Web Token alapú. A konfiguráció \texttt{Jwt} szakaszát a
\texttt{JwtOptions} beállításosztály (\texttt{Configuration/JwtOptions.cs}) képezi le, amely a
szimmetrikus aláírókulcsot, a kibocsátót (\texttt{HUNting.API}), a célközönséget
(\texttt{HUNting.Client}) és a token élettartamát (alapértelmezetten 60 perc) hordozza. A
bejövő kérések \texttt{Authorization} fejlécében érkező bearer-tokent a JWT-bearer middleware dolgozza
fel: ellenőrzi a kibocsátót, a célközönséget és az aláírást, a lejáratnál pedig mindössze
egyperces órakülönbséget tolerál. Sikeres validáció esetén a token claimjeiből előáll a kérés
identitása, amelyre a kontrollerek \texttt{[Authorize]} attribútumai és erőforrás-szintű
ellenőrzései épülnek. A tokenek \emph{kibocsátását} (a claimkészletet és az aláírást) a
felhasználókezelési fejezet (\ref{chap:felhasznalok}.~fejezet) mutatja be; ott esik szó a
megoldás korlátairól -- így a frissítő token és a tokenvisszavonás hiányáról -- is. Az
\texttt{appsettings.json}-ban tárolt aláírókulcs fejlesztői helykitöltő; éles környezetben
titokkezelő szolgáltatásból vagy környezeti változóból kellene érkeznie.

\subsection{Közös megvalósítási minták és a végpontkészlet}

Három megvalósítási minta a legtöbb kontrollerben visszatér, ezért itt kerül bemutatásra.

\paragraph{Validáció és hibakezelés.} Minden kontrollert az \texttt{[ApiController]} attribútum
jelöl, így a modellkötési és validációs hibákra a keretrendszer automatikusan 400-as választ ad.
A bemeneti szabályok első vonala az entitásokon elhelyezett DataAnnotations attribútumkészlet
(\texttt{[Required]}, \texttt{[MaxLength]} és társaik), amelyet az üzleti szabályok
kontrollerbeli, imperatív ellenőrzései egészítenek ki; a szolgáltatásrétegből érkező
\texttt{ArgumentException} jellegű hibákat a kontrollerek try/catch blokkjai fordítják 400-as
válaszra. Központi kivételkezelő middleware nincs: a nem kezelt kivétel a keretrendszer
alapértelmezett 500-as válaszát eredményezi.

\paragraph{Fájlfeltöltés.} A hirdetésképek, a társasági képek és a naplóképek feltöltése azonos
sémát követ: a végpontot a \texttt{[RequestSizeLimit(10\_485\_760)]} attribútum 10~MiB-ra
korlátozza; a kiterjesztés fehérlistán (\texttt{.jpg}, \texttt{.jpeg}, \texttt{.png},
\texttt{.webp}) ellenőrződik; a fájl ütközésmentes, GUID-alapú nevet kap, és a
\texttt{wwwroot/uploads} tartományspecifikus alkönyvtárába (\texttt{ads}, \texttt{societies},
\texttt{logs}) kerül; az adatbázisban csak a relatív URL tárolódik, amelyet a statikus
fájl-middleware szolgál ki; a rekord törlésekor a hozzá tartozó fájl is törlődik. A közös
lépéseket (kiterjesztés-ellenőrzés, könyvtár-előkészítés, mentés, törlés) a statikus
\texttt{UploadStorage} segédosztály (\texttt{Services/UploadStorage.cs}) végzi, így a
feltöltéskezelő kód nem ismétlődik kontrollerenként. A minta
korlátja, hogy a tartalom-ellenőrzés a kiterjesztésre szorítkozik (a fájl tényleges MIME-típusát
a szerver nem vizsgálja).

\paragraph{A végpontkészlet.} Az API nyolc kontrollere összesen 49 végpontot kínál;
a~\ref{tab:endpoint-osszes}.~táblázat ezek megoszlását összegzi. A részletes végponttáblák --
hozzáférési szabályokkal együtt -- az egyes funkciófejezetekben találhatók.

\begin{table}[htbp]
	\centering
	\caption{A HUNting.API kontrollerei és végpontjaik}
	\label{tab:endpoint-osszes}
	\begin{tabular}{|l|l|c|l|}
		\hline
		\textbf{Kontroller} & \textbf{Útvonalprefix} & \textbf{Végpont} & \textbf{Fejezet} \\
		\hline
		\texttt{AuthController} & \texttt{/api/Auth} & 4 & \ref{chap:felhasznalok}. \\
		\hline
		\texttt{UsersController} & \texttt{/api/Users} & 5 & \ref{chap:felhasznalok}. \\
		\hline
		\texttt{AdsController} & \texttt{/api/Ads} & 9 & \ref{chap:hirdetesek}. \\
		\hline
		\texttt{SocietiesController} & \texttt{/api/Societies} & 11 & \ref{chap:tarsasagok}. \\
		\hline
		\texttt{NewsController} & \texttt{/api/News} & 5 & \ref{chap:tarsasagok}. \\
		\hline
		\texttt{HuntingLogsController} & \texttt{/api/HuntingLogs} & 5 & \ref{chap:naplo}. \\
		\hline
		\texttt{MapController} & \texttt{/api/Map} & 9 & \ref{chap:terkep}. \\
		\hline
		\texttt{ChatController} & \texttt{/api/Chat} & 1 & \ref{chap:chatbot}. \\
		\hline
		\multicolumn{2}{|l|}{\textbf{Összesen}} & \textbf{49} & \\
		\hline
	\end{tabular}
\end{table}

\section{Frontend-alapok}
\label{sec:alapok-frontend}

\subsection{Belépési pont és provider-fa}

A kliensalkalmazás belépési pontja a \texttt{main.tsx}, amely a React 18 \texttt{createRoot}
API-jával, \texttt{StrictMode} burokban csatolja fel az \texttt{App} komponenst. Az \texttt{App}
(\texttt{App.tsx}) nem tartalmaz megjelenítési logikát: kizárólag az alkalmazást átfogó
szolgáltatók (providerek) fáját építi fel, amelyet a~\ref{fig:provider-fa}.~ábra szemléltet.
Legkívül a React Query \texttt{QueryClientProvider} áll, ezen belül az \texttt{AuthProvider}
biztosítja a hitelesítési állapotot, majd egy \texttt{Suspense} határ és a React Router
\texttt{RouterProvider} következik.

% TODO: FE-01 ábra (figures/fe-01-provider-fa.pdf) — fa-ábra: main → StrictMode → App →
% QueryClientProvider → AuthProvider → Suspense → RouterProvider → AppLayout/AuthLayout → oldalak.
\begin{figure}[htbp]
	\centering
	\fbox{\parbox[c][6cm][c]{0.85\textwidth}{\centering FE-01: A frontend provider- és
	layoutfája (elkészítendő)}}
	\caption{A kliensalkalmazás provider-hierarchiája és elrendezésrétegei}
	\label{fig:provider-fa}
\end{figure}

A szerverállapot kezelését a React Query végzi; a megosztott \texttt{queryClient} példány
(\texttt{lib/queryClient.ts}) három alapértelmezést rögzít: a lekérdezések eredménye egy percig
frissnek számít (\texttt{staleTime}), az ablakfókusz visszatérése nem vált ki újralekérést, a
sikertelen kérések pedig nem ismétlődnek automatikusan. A funkciómodulok konvenciója egységes: a
lekérdezések tartományspecifikus kulcsokkal (például \texttt{["ads", filters]}) regisztrálódnak,
a mutációk pedig sikeres lefutás után a \texttt{queryClient.invalidateQueries} hívással
érvénytelenítik az érintett kulcsok gyorsítótárát, így a felület automatikusan újratölti a friss
adatokat. Az \texttt{AuthProvider} és a mögötte álló \texttt{authStorage} modul működését -- a
token és a profil \texttt{localStorage}-beli tárolását, a munkamenet helyreállítását -- a
felhasználókezelési fejezet (\ref{chap:felhasznalok}.~fejezet) részletezi.

\subsection{Útvonalak és elrendezések}

Az útvonalkezelést a React Router \texttt{createBrowserRouter} hívása konfigurálja
(\texttt{routes/appRouter.tsx}), két elrendezés (layout) köré szervezve. Az \texttt{AppLayout}
a tartalmi oldalak kerete: rögzített fejlécet (\texttt{AppHeader}), legfeljebb \texttt{6xl}
szélességű tartalmi sávot, láblécet és -- minden oldalon elérhető, lebegő elemként -- a
chatbotot (\ref{chap:chatbot}.~fejezet) tartalmazza. Az \texttt{AuthLayout} a bejelentkezési és
regisztrációs oldal középre igazított kerete. A~\ref{tab:routes}.~táblázat a teljes
útvonaltáblát összegzi.

\begin{table}[htbp]
	\centering
	\caption{A kliensalkalmazás útvonalai}
	\label{tab:routes}
	\begin{tabular}{|l|l|l|}
		\hline
		\textbf{Útvonal} & \textbf{Oldal} & \textbf{Elrendezés} \\
		\hline
		\texttt{/} & \texttt{HomePage} & AppLayout \\
		\hline
		\texttt{/ads}, \texttt{/ads/new}, \texttt{/ads/:adId}, & hirdetésoldalak & AppLayout \\
		\texttt{/ads/:adId/edit} & & \\
		\hline
		\texttt{/societies}, \texttt{/societies/mine}, & társaságoldalak & AppLayout \\
		\texttt{/societies/manage}, \texttt{/societies/:societyId} & & \\
		\hline
		\texttt{/map} & \texttt{MapPage} & AppLayout \\
		\hline
		\texttt{/logs} & \texttt{HuntingLogsPage} & AppLayout \\
		\hline
		\texttt{/weather} & \texttt{WeatherPage} & AppLayout \\
		\hline
		\texttt{/profile} & \texttt{ProfilePage} & AppLayout \\
		\hline
		\texttt{/admin} & \texttt{AdminDashboardPage} & AppLayout \\
		\hline
		\texttt{/auth/login}, \texttt{/auth/register} & hitelesítési oldalak & AuthLayout \\
		\hline
		\texttt{*} & \texttt{NotFoundPage} & -- \\
		\hline
	\end{tabular}
\end{table}

Az ismeretlen útvonalakat egy globális, mindenre illeszkedő szabály és az \texttt{AppLayout}
\texttt{errorElement} beállítása egyaránt a \texttt{NotFoundPage} oldalra tereli. Két őszinte
megjegyzés kívánkozik ide. Egyrészt a router minden oldalkomponenst statikusan importál, azaz
\emph{nincs} kódszintű lusta betöltés (lazy loading): az \texttt{App} \texttt{Suspense} határa
jelenleg tartalék, egy későbbi kódfelbontás (code splitting) előkészítése, a teljes alkalmazás
egyetlen csomagban töltődik le. Másrészt az útvonaltábla nem tartalmaz központi útvonalőrt
(route guard): a védett oldalak hozzáférés-ellenőrzése oldalszinten, az egyes komponensekben
történik, ahogyan azt a felhasználókezelési fejezet bemutatja.

\subsection{API-kommunikáció}

A backenddel folytatott minden kommunikáció egyetlen megosztott axios-példányon, az
\texttt{apiClient}-en (\texttt{lib/apiClient.ts}) keresztül zajlik. Az alap-URL feloldása
háromlépcsős: ha a \texttt{VITE\_API\_BASE\_URL} környezeti változó be van állítva, az
érvényesül; fejlesztői módban az alapértelmezés a \texttt{https://localhost:5001/api} cím;
éles buildben a relatív \texttt{/api} útvonal. Megjegyzendő, hogy a backend fejlesztői profilja
valójában a 7003-as HTTPS-porton fut (\texttt{launchSettings.json}), így a fejlesztői
alapértelmezés önmagában nem talál célba -- a helyi fejlesztés a környezeti változó beállítását
igényli; ez a kódbázis ismert, javításra váró pontja. A példány viselkedésének magját a két
elfogó (interceptor) adja, amelyet a~\ref{lst:apiclient}.~kódrészlet mutat be.

\begin{lstlisting}[caption={Az \texttt{apiClient} kérés- és válasz-elfogója (kivonat)}, label={lst:apiclient}]
apiClient.interceptors.request.use((config) => {
  const token = getAccessToken();
  if (token) {
    config.headers.Authorization = `Bearer ${token}`;
  }
  return config;
});

apiClient.interceptors.response.use(
  (response) => response,
  (error) => {
    if (error.response?.status === 401) {
      purgeAuthStorage();
      window.location.href = "/auth/login";
    }
    if (error.response?.data?.message) {
      return Promise.reject(new Error(error.response.data.message));
    }
    // ... tovabbi hibauzenet-agak ...
  }
);
\end{lstlisting}

A kérés-elfogó a \texttt{localStorage}-ban tárolt hozzáférési tokent minden kimenő kérés
\texttt{Authorization} fejlécébe befűzi \texttt{Bearer} sémával, így az egyes funkciómodulok
API-hívásainak nem kell a hitelesítéssel foglalkozniuk. A válasz-elfogó két feladatot lát el:
a 401-es státuszú válaszokra globálisan reagál (törli a hitelesítési adatokat, és a
bejelentkezési oldalra irányít), minden más hibát pedig egységes \texttt{Error} objektummá
normalizál, a szerver \texttt{message} vagy \texttt{title} mezőjét részesítve előnyben -- a hívó
komponensek így mindig megjeleníthető hibaüzenetet kapnak. A modul harmadik közös eleme a
\texttt{resolveAssetUrl} segédfüggvény, amely a backend által visszaadott relatív képútvonalakat
(például \texttt{/uploads/ads/...}) teljes URL-lé egészíti ki az API gyökércíme alapján, az
abszolút (külső) URL-eket pedig változatlanul hagyja; ezt minden képet megjelenítő funkció
használja.

\subsection{Megjelenés és reszponzivitás}

A felület a Tailwind CSS 4-es verziójára épül, amelyben a téma nem külön konfigurációs fájlban,
hanem a globális stíluslapon (\texttt{index.css}), a \texttt{@theme} direktívával van definiálva.
A színvilág a vadászati tematikához igazodik: sötét erdőzöld alapszínek (háttér:
\texttt{\#111a13}, kártyák: \texttt{\#1f2a1f}) és arany-borostyán kiemelőszín
(\texttt{\#c3832f}) alkotják, a törzsszöveg az Inter, a címsorok a Playfair Display betűtípust
használják, a hátteret pedig rétegzett sugaras és lineáris színátmenet mélyíti. A reszponzivitás
mobil-első elvet követ, és a fejlécen figyelhető meg a legjobban: az \texttt{md} töréspont felett
vízszintes navigációs sáv és gombsor jelenik meg, alatta ezek elrejtőznek, és egy
hamburgermenü-gomb nyitja a lenyíló, érintésbarát navigációt. A fejléc rögzített, áttetsző,
elmosott hátterű sáv, a navigációs elemek pedig a bejelentkezett felhasználó fióktípusa és
szerepei szerint bővülnek (egyéni fióknál a „Társaságaim” és a „Naplók”, társasági fióknál a
„Társaságom”, adminisztrátornál az „Admin” menüpont). Megjegyzendő, hogy a felület magyar
feliratainak egy része jelenleg ékezet nélkül szerepel a forráskódban; ez ismert, javításra váró
szépséghiba.

% TODO: SC-01 képernyőkép (figures/sc-01-home-responsive.png) — főoldal desktop ÉS mobil
% nézetben egymás mellett; a mobil nézeten a hamburgermenü látszódjon.
\begin{figure}[htbp]
	\centering
	\fbox{\parbox[c][5.5cm][c]{0.85\textwidth}{\centering SC-01: A főoldal asztali és mobil
	nézetben (elkészítendő)}}
	\caption{A főoldal reszponzív megjelenése asztali és mobil nézetben}
	\label{fig:home-responsive}
\end{figure}

A fentiekben bemutatott alapréteg -- az egységes adatbázis-kontextus, a \texttt{Program.cs}-ben
összeállított backend-csővezeték és a provider-fára épülő kliensváz -- adja azt a közös nyelvet,
amelyre a következő fejezetek funkciónként építkeznek: mindegyikük az itt leírt mintákat
(JWT-validáció, \texttt{[Authorize]}, fájlfeltöltés, React Query-konvenciók,
\texttt{apiClient}) példányosítja a saját tartományára.

% !TeX root = ../main.tex
% 5. fejezet — Felhasználókezelés, jogosultságok és adminisztráció
% A ftsrg/thesis-template-latex sablonba \include-dal illesztendő tartalomfájl.
% Hivatkozott források: Backend/HUNting.API/Controllers/AuthController.cs, UsersController.cs,
%   Services/Auth/JwtTokenService.cs, Configuration/JwtOptions.cs, Services/UserService.cs,
%   Models/Users/ApplicationUser.cs, AccountType.cs, Program.cs, Data/SeedData.cs,
%   Frontend/HUNting/src/context/AuthContext.tsx, lib/authStorage.ts, lib/apiClient.ts,
%   services/authApi.ts, pages/Auth/LoginPage.tsx, RegisterPage.tsx, pages/Profile/ProfilePage.tsx,
%   pages/Admin/AdminDashboardPage.tsx, services/adminApi.ts

\chapter{Felhasználókezelés, jogosultságok és adminisztráció}
\label{chap:felhasznalok}

Ez a fejezet a dolgozat funkcióközpontú gerincének első tagja. A felhasználókezelés a HUNting
minden további funkciójának előfeltétele: a hirdetésfeladás, a társasági tagság, a vadásznapló és
a térképes rétegek egyaránt a bejelentkezett felhasználó személyazonosságára és szerepköreire
épülnek. A fejezet a bevezetőben rögzített egységes sablont követi: először a funkció célját és
felhasználói működését mutatom be, ezt követi az adatmodell, a backend- és a frontend-megvalósítás,
végül a teljes adatfolyam, valamint a biztonsági megfontolások és az ismert korlátok összegzése.

\section{A funkció célja és felhasználói működése}
\label{sec:users-cel}

A rendszer két, jogosultságaiban és felhasználói felületében is eltérő fióktípust különböztet meg,
amelyet a backend az \texttt{AccountType} felsorolt típussal ír le
(\texttt{Models/Users/AccountType.cs}): a \emph{Person} típus az egyéni vadászokat, a
\emph{Society} típus a vadásztársaságokat képviseli. Egyéni fiókkal a felhasználó hirdetést adhat
fel, társasági tagságot igényelhet, vadásznaplót vezethet és személyes térképrétegeket kezelhet.
A társasági fiók ezzel szemben a társaság adminisztrációs felületét éri el: tagfelvételi kérelmeket
bírál el, tagokat és bejegyzéseket kezel -- vadásznaplót azonban nem vezethet, és maga nem
jelentkezhet más társaságba.

A felhasználói folyamat a regisztrációval indul. A regisztrációs űrlapon egy kapcsolóval
választható a fióktípus; társasági fiók esetén az űrlap kiegészül a társaság adataival (név,
kapcsolattartási e-mail és telefon, rövid bemutatkozás). Sikeres regisztráció után a rendszer a
felhasználót azonnal be is lépteti: a válasz már tartalmazza a hozzáférési tokent, így külön
bejelentkezésre nincs szükség. A későbbi munkamenetek a bejelentkezési oldalon indulnak, ahol
e-mail-cím és jelszó megadásával igényelhető új token. A bejelentkezett felhasználó a
profiloldalon tekintheti át fiókadatait és szerepköreit, innen érhetők el a fióktípushoz kötött
felületek (egyéni fióknál a „Társaságaim”, társasági fióknál a „Társaság kezelése” oldal), és itt
található a kijelentkezés is. Védett oldal vendégként történő megnyitásakor az alkalmazás a
bejelentkezési oldalra irányít át.

% TODO: SC-12 képernyőkép (figures/sc-12-login.png) — a LoginPage mojibake-javítása UTÁN készüljön!
\begin{figure}[htbp]
	\centering
	\fbox{\parbox[c][5.5cm][c]{0.85\textwidth}{\centering SC-12: A bejelentkezési oldal képernyőképe (elkészítendő)}}
	\caption{A HUNting bejelentkezési oldala}
	\label{fig:login-screen}
\end{figure}

% TODO: SC-13 képernyőkép (figures/sc-13-register.png) — a fióktípus-váltó és a társasági mezők látszódjanak.
\begin{figure}[htbp]
	\centering
	\fbox{\parbox[c][5.5cm][c]{0.85\textwidth}{\centering SC-13: A regisztrációs oldal társasági mezőkkel (elkészítendő)}}
	\caption{Regisztráció társasági fióktípussal, a kiegészítő társasági adatokkal}
	\label{fig:register-screen}
\end{figure}

\section{Adatmodell}
\label{sec:users-adatmodell}

A felhasználókat az ASP.NET Core Identity keretrendszer kezeli, amelyet a projekt egész számú
elsődleges kulcsokkal használ. Az alkalmazás saját felhasználó-entitása, az
\texttt{ApplicationUser} az \texttt{IdentityUser<int>} osztályból származik, és négy
tartományspecifikus mezővel egészíti ki azt: a kötelező, legfeljebb 100 karakteres
\texttt{FullName} névmezővel, az opcionális \texttt{Phone} telefonszámmal, a létrehozás
időbélyegével (\texttt{CreatedAt}), valamint a fióktípust tároló \texttt{AccountType} mezővel
(\texttt{Models/Users/ApplicationUser.cs}). Az entitás navigációs tulajdonságai a felhasználó
kapcsolatait írják le: az \texttt{OwnedSociety} az általa birtokolt társaságra mutat (egy-egy
kapcsolat), a \texttt{Societies} kollekció a tagságait, a \texttt{MembershipRequests} a benyújtott
tagfelvételi kérelmeit, a \texttt{HuntingLogs} pedig a vadásznapló-bejegyzéseit tartalmazza.

Az Identity alapértelmezett \texttt{AspNet*} előtagú tábláit a \texttt{HuntingContext} saját,
beszédesebb nevekre képezi le (\texttt{Users}, \texttt{Roles}, \texttt{UserRoles},
\texttt{UserClaims}, \texttt{UserLogins}, \texttt{RoleClaims}, \texttt{UserTokens}); az átnevezést
a \texttt{RenameIdentityTables} migráció vezette át az adatbázissémán. Az e-mail-cím egyedisége
két szinten is biztosított: az Identity konfigurációjában a \texttt{RequireUniqueEmail} beállítás,
az adatbázisban pedig egyedi index garantálja. A társaság és tulajdonosa közötti egy-egy
kapcsolatot a \texttt{Society.OwnerUserId} idegen kulcs valósítja meg; a tulajdonos felhasználó
törlésekor a kapcsolat \texttt{SetNull} viselkedéssel szakad el, így a társaság rekordja nem vész
el. Megjegyzendő, hogy a \texttt{Society} entitás saját \texttt{Email} és \texttt{PasswordHash}
mezőkkel is rendelkezik, a bejelentkezés azonban kizárólag a felhasználói fiókon keresztül
történik: társasági regisztrációkor a \texttt{PasswordHash} mezőbe csupán egy véletlen generált
helykitöltő érték kerül.

Az alkalmazásszintű szerepköröket az Identity szereptáblája tárolja. A kezdeti adatfeltöltés
(\texttt{Data/SeedData.cs}) három szerepet hoz létre: az \emph{Administrator} a teljes rendszert
felügyeli, a \emph{SocietyManager} a társasági fiókok alapértelmezett szerepe, a \emph{Hunter}
pedig az egyéni vadászoké. Fontos megkülönböztetni ezektől a társaságon belüli szerepeket
(\texttt{SocietyRole}: Vezeto, Gazda, Vador, Tag), amelyek nem az Identity részei, hanem a tagsági
kapcsolatrekordhoz kötődnek; ezeket a~\ref{chap:tarsasagok}.~fejezet tárgyalja.

\section{Backend-megvalósítás}
\label{sec:users-backend}

A hitelesítési funkciókat az \texttt{AuthController} valósítja meg, amely -- a rétegzett
felépítéstől némileg eltérve -- nem egy dedikált szolgáltatásosztályra, hanem közvetlenül az
Identity \texttt{UserManager} és \texttt{SignInManager} komponenseire, a tokengeneráláshoz pedig
az \texttt{ITokenService} absztrakcióra támaszkodik. A~\ref{tab:auth-endpoints}.~táblázat a
kontroller végpontjait foglalja össze.

\begin{table}[htbp]
	\centering
	\caption{Az \texttt{AuthController} végpontjai}
	\label{tab:auth-endpoints}
	\begin{tabular}{|l|l|p{6.6cm}|}
		\hline
		\textbf{Végpont} & \textbf{Hozzáférés} & \textbf{Feladat} \\
		\hline
		POST \texttt{/api/Auth/register} & anonim & Fiók létrehozása (Person vagy Society), azonnali tokenkibocsátással \\
		\hline
		POST \texttt{/api/Auth/login} & anonim & Hitelesítés, JWT és felhasználói adatok visszaadása \\
		\hline
		POST \texttt{/api/Auth/logout} & bejelentkezett & Kijelentkezés (állapotmentes, 204-es válasz) \\
		\hline
		GET \texttt{/api/Auth/me} & bejelentkezett & Az aktuális felhasználó profilja, szerepei és birtokolt társasága \\
		\hline
	\end{tabular}
\end{table}

\subsection{Regisztráció}

A \texttt{Register} akció a kérés DTO-jából (\texttt{RegisterRequest} rekord) felépíti az
\texttt{ApplicationUser} példányt, majd a \texttt{UserManager.CreateAsync} hívással hozza létre a
fiókot; a jelszó hash-eléséről az Identity gondoskodik. Sikertelen létrehozás esetén a végpont
szándékosan általános hibaüzenettel tér vissza, a konkrét okok felsorolása nélkül. Ezt követően a
fiók a típusának megfelelő alapértelmezett szerepet kapja: társasági fiók esetén
\emph{SocietyManager}, egyébként \emph{Hunter} szerepet.

Társasági regisztrációnál a folyamat két további lépéssel bővül. A kontroller ellenőrzi, hogy a
kérés tartalmazza-e a társaság nevét -- ennek hiányában hibával tér vissza --, majd létrehozza a
\texttt{Society} entitást, tulajdonosként az imént regisztrált felhasználót megjelölve. Ha a
seed-adatok között létezik a „Vezeto” elnevezésű társasági szerep, a rendszer a tulajdonost azonnal
taggá is teszi ezzel a szereppel, így a társasági fiók a regisztráció pillanatától teljes körűen
kezelheti saját társaságát. A válasz mindkét ágon a bejelentkezésével azonos szerkezetű: a
kibocsátott tokent, annak lejáratát és a felhasználó adatait (\texttt{UserDto}) tartalmazza.

\subsection{Bejelentkezés és tokenkibocsátás}

A \texttt{Login} akció először e-mail-cím alapján keresi meg a felhasználót
(\texttt{FindByEmailAsync}), majd a \texttt{SignInManager.CheckPasswordSignInAsync} hívással
ellenőrzi a jelszót. Mindkét hibaág ugyanazt az általános üzenetet adja vissza („Hibás e-mail cím
vagy jelszó.”), így a válaszból nem derül ki, hogy létezik-e fiók az adott címmel -- ez a
felhasználó-felsorolásos (user enumeration) támadások elleni bevett védekezés. Sikeres hitelesítés
után a \texttt{JwtTokenService} állítja elő a hozzáférési tokent
(\ref{lst:jwt-token}.~kódrészlet).

\begin{lstlisting}[language={[Sharp]C}, caption={A JWT előállítása a \texttt{JwtTokenService} osztályban (kivonat)}, label={lst:jwt-token}]
var claims = new List<Claim>
{
    new(JwtRegisteredClaimNames.Sub, user.Id.ToString()),
    new(JwtRegisteredClaimNames.Email, user.Email ?? string.Empty),
    new(JwtRegisteredClaimNames.Jti, Guid.NewGuid().ToString()),
    new(ClaimTypes.NameIdentifier, user.Id.ToString()),
    new(ClaimTypes.Name, user.FullName),
    new(ClaimTypes.Email, user.Email ?? string.Empty)
};

foreach (var role in roles)
{
    claims.Add(new Claim(ClaimTypes.Role, role));
}

claims.Add(new("account_type", user.AccountType.ToString()));

var expires = DateTime.UtcNow.AddMinutes(_options.AccessTokenLifetimeMinutes);

var jwt = new JwtSecurityToken(
    issuer: _options.Issuer,
    audience: _options.Audience,
    claims: claims,
    expires: expires,
    signingCredentials: credentials);
\end{lstlisting}

A token a felhasználó azonosítóját, nevét és e-mail-címét, az összes alkalmazásszintű szerepét,
valamint egy egyedi tokenazonosítót (\texttt{jti}) és a fióktípust leíró saját
\texttt{account\_type} claimet hordozza. Az aláírás szimmetrikus kulcsú HMAC-SHA256 algoritmussal
készül; a kulcsot, a kibocsátót, a célközönséget és az élettartamot (alapértelmezetten 60 perc) a
konfiguráció \texttt{Jwt} szakaszából kötött \texttt{JwtOptions} beállításosztály szolgáltatja. A
bejövő kérések tokenjeit a~\ref{chap:alapok}.~fejezetben bemutatott bearer-middleware validálja
(kibocsátó-, célközönség- és aláírásellenőrzés, egyperces megengedett órakülönbséggel). A
\texttt{Me} végpont a validált tokenből a \texttt{NameIdentifier} claim alapján azonosítja a
felhasználót, és a friss adatbázis-állapotot adja vissza. A \texttt{Logout} végpont állapotmentes:
kizárólag 204-es választ ad, a token érvényességét nem vonja vissza -- ennek következményeit
a~\ref{sec:users-biztonsag}.~alfejezet tárgyalja.

\subsection{A jogosultsági modell}

A hozzáférés-szabályozás két szinten valósul meg. Az első szint deklaratív: a végpontokat az
ASP.NET Core \texttt{[Authorize]} attribútuma védi, amely szerepkör-megkötéssel is ellátható -- a
felhasználó-adminisztrációs kontroller például teljes egészében az
\texttt{[Authorize(Roles = "Administrator")]} attribútummal zárja ki a nem adminisztrátor
hívókat. A második szint imperatív, erőforrás-szintű ellenőrzés: a módosító műveletek a kérdéses
rekord tulajdonosát vetik össze a tokenből kiolvasott felhasználó-azonosítóval, és az
adminisztrátori szerep ez alól kivételt képez. Ezt a mintát szemlélteti
a~\ref{lst:owner-check}.~kódrészlet a hirdetésmódosítás példáján; ugyanez az ellenőrzés ismétlődik
a hirdetéstörlésben, a képkezelő végpontokban, a vadásznaplóban és a társasági bejegyzéseknél is.

\begin{lstlisting}[language={[Sharp]C}, caption={Erőforrás-tulajdonos ellenőrzés az \texttt{AdsController.Update} akcióban (kivonat)}, label={lst:owner-check}]
var existing = await _svc.GetByIdAsync(id);
if (existing == null) return NotFound();

var userId = GetUserId();
if (userId == null) return Unauthorized();

var isAdmin = User.IsInRole("Administrator");
if (!isAdmin && existing.OwnerUserId != userId.Value)
{
    return Forbid();
}

// ... mezok masolasa ...

if (isAdmin)
{
    existing.OwnerUserId = ad.OwnerUserId;
    existing.OwnerSocietyId = ad.OwnerSocietyId;
}
\end{lstlisting}

A jogosultsági modell harmadik eleme a fióktípus szerinti korlátozás: egyes végpontok a token
\texttt{account\_type} claimje, illetve a felhasználó rekordja alapján zárják ki a társasági
fiókokat -- társasági fiók nem hozhat létre vadásznapló-bejegyzést
(\ref{chap:naplo}.~fejezet), és nem nyújthat be tagfelvételi kérelmet
(\ref{chap:tarsasagok}.~fejezet).

\section{Frontend-megvalósítás}
\label{sec:users-frontend}

A kliensoldali hitelesítés központi eleme az \texttt{AuthContext} React-kontextus
(\texttt{src/context/AuthContext.tsx}), amely az egész alkalmazást átfogó szolgáltatóként a
bejelentkezett felhasználót és a hitelesítési műveleteket (\texttt{login}, \texttt{register},
\texttt{logout}, \texttt{refreshProfile}) teszi elérhetővé. A munkamenet két értéken alapul,
amelyeket az \texttt{authStorage} segédmodul a böngésző \texttt{localStorage} tárolójában kezel: a
\texttt{hunting.accessToken} kulcs a JWT-t, a \texttt{hunting.user} kulcs a felhasználó
szerializált profilját tárolja. A kontextus kezdeti állapotát a tárolt profilból állítja helyre,
így az oldal újratöltése nem szakítja meg a munkamenetet; ha tokent talál profil nélkül, a
felcsatolódáskor a \texttt{/auth/me} végpontról frissíti azt, sikertelenség esetén pedig kiüríti a
tárolót.

A bejelentkezési és regisztrációs művelet azonos sémát követ: az \texttt{authApi} szolgáltatásmodul
meghívja a megfelelő végpontot -- a regisztrációnál a kliensoldali \texttt{name} mezőt a backend
által várt \texttt{fullName} mezőre képezve át --, majd sikeres válasz esetén a token és a profil a
tárolóba és a kontextus állapotába kerül. A hibákat a hívó űrlap kapja meg és jeleníti meg. A
kijelentkezés először a backend \texttt{logout} végpontját hívja (ha van tárolt token), majd minden
esetben kiüríti a tárolót és a kontextust.

A kéréseket az \texttt{apiClient} axios-példány két elfogója (interceptor) egészíti ki. A
kérés-elfogó minden hívás \texttt{Authorization} fejlécébe befűzi a tárolt tokent
\texttt{Bearer} sémával. A válasz-elfogó a 401-es státuszú válaszokra globálisan reagál: törli a
hitelesítési kulcsokat és a bejelentkezési oldalra irányítja a böngészőt -- így a lejárt vagy
érvénytelen tokennel dolgozó munkamenet bármely védett hívás hatására automatikusan és egységesen
kijelentkezik.

Az oldalszintű hozzáférés-védelem nem központi útvonalőrrel (route guard), hanem az egyes
oldalkomponensekben valósul meg. A profiloldal vendég látogatóját a \texttt{Navigate} komponens a
bejelentkezési oldalra irányítja; a bejelentkezési és regisztrációs oldal a már hitelesített
felhasználót az útvonalállapotban (\texttt{location.state.from}) átadott eredeti céloldalra --
ennek hiányában a főoldalra -- lépteti tovább; az admin vezérlőpult pedig a szerepkörlista alapján
dönt. A regisztrációs űrlap kliensoldali validációt is végez: a kötelező mezőket és az e-mail
formátumát natív HTML-attribútumok ellenőrzik, a jelszóra hatkarakteres minimum vonatkozik
(összhangban a szerveroldali Identity-szabályzattal), társasági fióktípusnál pedig a beküldés
előfeltétele a társaságnév kitöltése.

% TODO: SC-16 képernyőkép (figures/sc-16-profile.png)
\begin{figure}[htbp]
	\centering
	\fbox{\parbox[c][5.5cm][c]{0.85\textwidth}{\centering SC-16: A profiloldal képernyőképe (elkészítendő)}}
	\caption{A profiloldal a fiókadatokkal, szerepkörökkel és a fióktípusfüggő gyorslinkekkel}
	\label{fig:profile-screen}
\end{figure}

\section{Adminisztrációs felület}
\label{sec:users-admin}

A rendszergazdai funkciók két rétegből állnak. A backendoldali \texttt{UsersController} teljes
felhasználó-adminisztrációt kínál (listázás, lekérdezés, létrehozás, módosítás, törlés), és --
mint fentebb szerepelt -- kizárólag \emph{Administrator} szereppel érhető el. A létrehozó és
módosító végpontok a szerepkiosztást is kezelik: a \texttt{UserService} a még nem létező szerepet
automatikusan megteremti, a módosító akció pedig halmazkülönbség-képzéssel számítja ki az
eltávolítandó és hozzáadandó szerepeket, így a kérésben megadott lista deklaratívan írja le a
kívánt végállapotot.

A frontendoldali admin vezérlőpult (\texttt{AdminDashboardPage}) két nézet között vált:
hirdetések és felhasználók. Mindkét nézet lusta lekérdezést használ -- a React Query
\texttt{enabled} kapcsolójával csak az aktív nézet adatai töltődnek be --, a törlőműveletek pedig
natív megerősítő kérdés (\texttt{window.confirm}) után futnak le, majd a találati lista
gyorsítótárát érvénytelenítve azonnal frissítik a felületet. A felület jelenleg listázásra és
törlésre szorítkozik; a backend létrehozó és módosító végpontjaihoz nem tartozik felületi elem,
ami a~\ref{sec:users-biztonsag}.~alfejezetben felsorolt korlátok egyike. Az adminisztrátor ezen
felül a vadásznaplók teljes állományát is lekérdezheti a \texttt{scope=all} paraméterrel
(\ref{chap:naplo}.~fejezet).

% TODO: SC-14 képernyőkép (figures/sc-14-admin.png) — felhasználók nézet, törlés-megerősítéssel.
\begin{figure}[htbp]
	\centering
	\fbox{\parbox[c][5.5cm][c]{0.85\textwidth}{\centering SC-14: Az admin vezérlőpult képernyőképe (elkészítendő)}}
	\caption{Az adminisztrációs vezérlőpult felhasználókezelő nézete}
	\label{fig:admin-screen}
\end{figure}

\section{Adatfolyam, biztonság és korlátok}
\label{sec:users-biztonsag}

A bejelentkezés teljes adatfolyama a következő lépésekből áll (a folyamatot
a~\ref{fig:auth-sequence}.~ábra szemlélteti):

\begin{enumerate}
	\item A felhasználó kitölti a bejelentkezési űrlapot; a \texttt{LoginPage} az
	      \texttt{AuthContext.login} műveletet hívja.
	\item Az \texttt{authApi} a \texttt{POST /api/auth/login} végpontra küldi az e-mail-címet és a
	      jelszót; a kérés-elfogó ekkor még nem fűz be tokent.
	\item Az \texttt{AuthController} a \texttt{UserManager} segítségével megkeresi a felhasználót,
	      majd a \texttt{SignInManager} ellenőrzi a jelszót a tárolt hash-sel szemben.
	\item Siker esetén a \texttt{JwtTokenService} előállítja az aláírt tokent a szerepkör- és
	      identitás-claimekkel.
	\item A válaszból a kliens a tokent és a profilt a \texttt{localStorage}-ba menti, és frissíti a
	      kontextus állapotát; a felület a bejelentkezett nézetre vált.
	\item A további API-hívások \texttt{Authorization: Bearer} fejléccel utaznak; a middleware
	      kérésenként validálja a tokent, és a claimekből állítja elő a felhasználó identitását.
	\item Érvénytelen vagy lejárt token esetén a backend 401-gyel válaszol, amire a válasz-elfogó a
	      tároló kiürítésével és a bejelentkezési oldalra történő átirányítással reagál.
\end{enumerate}

% TODO: B-01 szekvenciadiagram (figures/b-01-login-sequence.pdf) — LoginPage → AuthContext → axios →
%       AuthController → SignInManager → JwtTokenService → localStorage.
\begin{figure}[htbp]
	\centering
	\fbox{\parbox[c][6.5cm][c]{0.9\textwidth}{\centering B-01: A bejelentkezés és tokenkibocsátás szekvenciadiagramja (elkészítendő)}}
	\caption{A bejelentkezési folyamat szekvenciadiagramja}
	\label{fig:auth-sequence}
\end{figure}

% TODO: B-02 folyamatábra (figures/b-02-request-auth.pdf) — kérések hitelesítése: Bearer-fejléc →
%       validálás → claims → [Authorize]/szerep → erőforrás-tulajdonos ellenőrzés; 401-interceptor ág.
\begin{figure}[htbp]
	\centering
	\fbox{\parbox[c][6.5cm][c]{0.9\textwidth}{\centering B-02: A kérések hitelesítésének folyamatábrája (elkészítendő)}}
	\caption{A védett kérések kiszolgálásának folyamata a tokenvalidálástól az erőforrás-szintű ellenőrzésig}
	\label{fig:auth-flow}
\end{figure}

A megoldás több ponton követi a bevett biztonsági gyakorlatot: a jelszavakat az Identity
alapértelmezett, PBKDF2 alapú hash-elése védi; a bejelentkezési hibaüzenet nem árulja el a fiók
létezését; az e-mail-cím egyedisége adatbázisszinten kikényszerített; a token rövid élettartamú és
aláírása kérésenként ellenőrzött; az erőforrás-szintű tulajdonos-ellenőrzés pedig megakadályozza,
hogy egy hitelesített felhasználó más rekordjait módosítsa.

Ugyanakkor a jelenlegi megvalósítás korlátait is rögzíteni kell. A jelszószabályzat tudatosan
enyhített -- hat karakteres minimum, összetettségi követelmények nélkül --, ami éles üzemben
szigorítandó. A kijelentkezés kizárólag kliensoldali: a kibocsátott token a lejáratáig érvényes
marad, mivel sem visszavonási lista, sem frissítőtoken-mechanizmus nincs. A token
\texttt{localStorage}-ban tárolása egyszerű és újratöltésálló, de XSS-támadás esetén a token
kiolvasható -- ezt a kockázatot a keretrendszer kimeneti kódolása mérsékli, meg nem szünteti. A
társasági regisztráció nem atomi művelet: ha a felhasználó létrehozása után a társaságnév
hiányzik, a kérés hibával tér vissza, miközben a felhasználói fiók már létrejött. Az admin
felhasználó-létrehozó és -módosító végpontok kivételei kezeletlenül 500-as válasszá alakulnak,
felületi funkció pedig nem tartozik hozzájuk. Végül a kezdeti adatfeltöltésben rögzített
adminisztrátori jelszó kizárólag fejlesztési környezetben elfogadható; éles telepítés előtt
cserélendő, ahogyan a konfigurációban tárolt JWT-aláírókulcs is titokkezelőbe való.

% !TeX root = ../main.tex
% 6. fejezet — Hirdetési piactér
% Hivatkozott források: Backend/HUNting.API/Controllers/AdsController.cs, Services/AdService.cs,
%   Services/AdQueryOptions.cs, Models/PagedResult.cs, Models/Ads/*, Data/HuntingContext.cs,
%   Frontend/HUNting/src/services/adsApi.ts, pages/Ads/AdsListPage.tsx, AdCreatePage.tsx,
%   AdEditPage.tsx, AdDetailPage.tsx, AdForm.tsx

\chapter{Hirdetési piactér}
\label{chap:hirdetesek}

A hirdetési piactér a HUNting legnagyobb forgalmú, publikusan is elérhető funkcióterülete:
vadászati eszközök, felszerelések és vendégvadász-meghívások adásvételét támogatja. A fejezet a
funkciót a már ismertetett sablon szerint, az adatmodelltől a felhasználói felületig mutatja be,
külön kitérve a szerveroldali lapozás-szűrés és a képfeltöltés megvalósítására.

\section{A funkció célja és felhasználói működése}
\label{sec:ads-cel}

A piactér böngészése nem igényel bejelentkezést: a hirdetéslista kulcsszó szerint kereshető, és
kategória, valamint minimum- és maximumár szerint szűrhető; a találatok laponként kilencesével,
létrehozási idő szerint csökkenő sorrendben jelennek meg. A listaelemről a hirdetés részletező
oldalára lehet továbblépni, amely a leírást, az árat, a helyszínt, a feltöltött képeket és a
kapcsolatfelvételi adatokat mutatja. Bejelentkezett felhasználó új hirdetést adhat fel: az űrlapon
cím, leírás, kategória, ár, helyszín és elérhetőségek adhatók meg, a képek pedig húzd-és-ejtsd
módszerrel, vágólapról beillesztve vagy fájlböngészővel csatolhatók. A hirdetés tulajdonosa -- és
az adminisztrátor -- a részletező oldalról szerkesztheti és törölheti a hirdetést, valamint
egyenként távolíthatja el a képeket; bejelentkezett felhasználó a listaoldalon a „Saját
hirdetéseim” kapcsolóval csak a maga hirdetéseire szűkítheti a találatokat.

% TODO: SC-02 képernyőkép (figures/sc-02-ads-list.png) — aktív kategória- és árszűrővel.
\begin{figure}[htbp]
	\centering
	\fbox{\parbox[c][5cm][c]{0.85\textwidth}{\centering SC-02: A hirdetéslista szűrőkkel és lapozóval (elkészítendő)}}
	\caption{A hirdetési piactér listaoldala keresővel, kategória- és árszűrőkkel}
	\label{fig:ads-list-screen}
\end{figure}

% TODO: SC-03 képernyőkép (figures/sc-03-ad-detail.png) — tulajdonosként, a kezelőgombokkal.
\begin{figure}[htbp]
	\centering
	\fbox{\parbox[c][5cm][c]{0.85\textwidth}{\centering SC-03: A hirdetés részletező oldala (elkészítendő)}}
	\caption{Hirdetés részletei képekkel, kapcsolati adatokkal és tulajdonosi műveletekkel}
	\label{fig:ad-detail-screen}
\end{figure}

\section{Adatmodell}
\label{sec:ads-adatmodell}

A funkció négy entitásra épül (\texttt{Models/Ads/}). A központi \texttt{Ad} entitás tárolja a
hirdetés címét, leírását, árát (\texttt{decimal?}), a helyszín szöveges megnevezését és opcionális
koordinátáit, a kapcsolatfelvételi e-mail-címet és telefonszámot, az állapotkódot, valamint a
létrehozás időbélyegét. A kategóriákat az \texttt{AdCategory} szótártábla írja le, amelynek neve
egyedi indexszel védett; a kezdeti adatfeltöltés tizenkét kategóriát hoz létre (például
„Fegyverek”, „Lőszerek”, „Vadkamera”). A hirdetés állapotát az \texttt{AdStatus} referenciatábla
kódolja -- itt nem numerikus, hanem karakterláncos elsődleges kulcsot használ a rendszer
(\texttt{ACTIVE}, \texttt{PENDING}, \texttt{SOLD}, \texttt{HIDDEN}) --, így az állapotkód
önmagában is beszédes. A képeket az \texttt{AdImage} entitás kapcsolja a hirdetéshez: a tárolt
útvonal mellett \texttt{SortOrder} rendezőszámot is tartalmaz, és a hirdetés törlésekor
kaszkádolva törlődik.

A tulajdonlás a modell legérdekesebb megszorítása. A hirdetés gazdája kétféle lehet --
felhasználó (\texttt{OwnerUserId}) vagy társaság (\texttt{OwnerSocietyId}) --, és a
\texttt{HuntingContext} egy adatbázisszintű ellenőrző megszorítással
(\texttt{CK\_Ads\_SingleOwner}) garantálja, hogy a két idegen kulcs közül \emph{pontosan az egyik}
legyen kitöltve. Ugyanezt az invariánst az üzleti réteg is ellenőrzi (lásd
a~\ref{sec:ads-backend}.~alfejezetet), így a szabály megsértése már az adatbázis elérése előtt
kiderül. Mindkét tulajdonosi idegen kulcs kaszkád törlésű: a felhasználó vagy a társaság
törlésekor a hirdetései is törlődnek; a kategória- és státuszhivatkozás ezzel szemben
\texttt{Restrict} viselkedésű, tehát használatban lévő szótárérték nem törölhető.

\section{Backend-megvalósítás}
\label{sec:ads-backend}

A funkciót az \texttt{AdsController} és a mögötte álló \texttt{AdService} szolgáltatás valósítja
meg; a kontroller kizárólag a HTTP-réteggel és a jogosultság-ellenőrzéssel foglalkozik, a
lekérdezés-összeállítás és az adatműveletek a szolgáltatásban élnek. A végpontokat
a~\ref{tab:ads-endpoints}.~táblázat foglalja össze.

\begin{table}[htbp]
	\centering
	\caption{Az \texttt{AdsController} végpontjai}
	\label{tab:ads-endpoints}
	\begin{tabular}{|l|l|p{5.9cm}|}
		\hline
		\textbf{Végpont} & \textbf{Hozzáférés} & \textbf{Feladat} \\
		\hline
		GET \texttt{/api/Ads} & anonim & Lapozott, szűrhető hirdetéslista \\
		\hline
		GET \texttt{/api/Ads/categories} & anonim & Kategórialista \\
		\hline
		GET \texttt{/api/Ads/statuses} & anonim & Állapotlista \\
		\hline
		GET \texttt{/api/Ads/\{id\}} & anonim & Egy hirdetés részletei \\
		\hline
		POST \texttt{/api/Ads} & bejelentkezett & Hirdetés létrehozása \\
		\hline
		PUT \texttt{/api/Ads/\{id\}} & tulajdonos v. admin & Hirdetés módosítása \\
		\hline
		DELETE \texttt{/api/Ads/\{id\}} & tulajdonos v. admin & Hirdetés törlése \\
		\hline
		POST \texttt{/api/Ads/\{id\}/images} & tulajdonos v. admin & Kép feltöltése (max.\ 10~MiB) \\
		\hline
		DELETE \texttt{/api/Ads/\{id\}/images/\{imageId\}} & tulajdonos v. admin & Kép törlése \\
		\hline
	\end{tabular}
\end{table}

\subsection{Lapozott és szűrt lekérdezés}

A listázó végpont a lekérdezési paramétereket az \texttt{AdQueryOptions} osztályba köti
(oldalszám, lapméret, keresőkifejezés, kategória, állapotkód, ár-intervallum, tulajdonos). A
szolgáltatás a lapozási értékeket normalizálja -- egynél kisebb oldalszámból egy lesz, az 1 és 100
közé nem eső lapméret helyett az alapértelmezett tizenkettő érvényesül --, majd a szűrőfeltételeket
csak akkor fűzi a lekérdezéshez, ha az adott paraméter ténylegesen meg van adva
(\ref{lst:ads-paged}.~kódrészlet). Így minden szűrés az adatbázisban, SQL-szinten fut le, a
válasz pedig a \texttt{PagedResult<T>} rekordban a találatok mellett a teljes találatszámot és a
lapozási adatokat is visszaadja -- a kliens ebből számítja az oldalszámot.

\begin{lstlisting}[language={[Sharp]C}, caption={Szerveroldali lapozás és feltételes szűrés az \texttt{AdService.GetPagedAsync} metódusban (kivonat)}, label={lst:ads-paged}]
var page = options.Page < 1 ? 1 : options.Page;
var take = options.Take is < 1 or > 100 ? 12 : options.Take;

IQueryable<Ad> query = _context.Ads
    .AsNoTracking()
    .Include(a => a.Category)
    .Include(a => a.Images)
    .Include(a => a.Status);

if (!string.IsNullOrWhiteSpace(options.Search))
{
    var term = options.Search.Trim();
    query = query.Where(a =>
        EF.Functions.Like(a.Title, $"%{term}%") ||
        (a.Description != null && EF.Functions.Like(a.Description, $"%{term}%")));
}

if (options.MinPrice.HasValue)
{
    query = query.Where(a => a.Price >= options.MinPrice.Value);
}

var total = await query.CountAsync();

var items = await query
    .OrderByDescending(a => a.CreatedAt)
    .Skip((page - 1) * take)
    .Take(take)
    .ToListAsync();

return new PagedResult<Ad>(items, total, page, take);
\end{lstlisting}

A keresés az \texttt{EF.Functions.Like} függvénnyel fordul SQL \texttt{LIKE} kifejezésre, és a
címre, valamint a leírásra egyaránt illeszkedik. A lekérdezés \texttt{AsNoTracking} módban fut,
mivel az eredmény csak olvasásra szolgál, a kapcsolódó kategória-, kép-, állapot- és
tulajdonos-adatokat pedig \texttt{Include} hívások töltik be egyetlen lekérdezésben.

\subsection{Írási műveletek és tulajdonlás}

Hirdetés létrehozásakor a kontroller a hívó szerepétől függően kezeli a tulajdonlást: nem
adminisztrátor esetén a tulajdonos kényszerítetten maga a hívó (a kérésben esetleg megadott
tulajdonosmezők felülíródnak), adminisztrátor viszont tetszőleges felhasználói vagy társasági
tulajdonost adhat meg -- ha egyiket sem adja meg, a rendszer őt magát jelöli tulajdonosnak. Az
\texttt{AdService} minden írási művelet előtt a „pontosan egy tulajdonos” invariánst is ellenőrzi,
és megsértése esetén \texttt{ArgumentException} kivételt dob, amelyet a kontroller 400-as
válaszra fordít. A módosítás és a törlés a~\ref{chap:felhasznalok}.~fejezetben bemutatott
erőforrás-tulajdonos mintát követi (\ref{lst:owner-check}.~kódrészlet): a műveletet csak a
tulajdonos vagy adminisztrátor végezheti el, és a tulajdonosmezők átírása adminisztrátori
kiváltság.

\subsection{Képfeltöltés}

A képfeltöltő végpont többrészes (\texttt{multipart/form-data}) kérést fogad, amelynek méretét a
\texttt{[RequestSizeLimit]} attribútum 10~MiB-ban korlátozza. A feldolgozás több lépcsős
védelmet alkalmaz: üres fájl esetén azonnali hibával tér vissza, a fájlnév kiterjesztését
kisbetűsítve fehérlistával veti össze (csak \texttt{.jpg}, \texttt{.jpeg}, \texttt{.png} és
\texttt{.webp} engedélyezett), a lemezre pedig nem az eredeti fájlnéven, hanem újonnan generált
GUID néven ír -- ezzel kizárva az útvonal-befecskendezést és a fájlnév-ütközést. A fájl a webes
gyökér \texttt{uploads/ads} könyvtárába kerül (a könyvtárat a kontroller szükség esetén
létrehozza), az adatbázisba csak a relatív URL íródik, a kiszolgálást pedig a statikus
fájl-middleware végzi. Az \texttt{AdService.AddImageAsync} a rendezőszámot a meglévő képek
maximuma alapján, folytatólagosan osztja ki. Kép törlésekor a végpont ellenőrzi, hogy a kép
valóban az útvonalban megadott hirdetéshez tartozik-e, majd a lemezen tárolt fájlt és az
adatbázisrekordot is eltávolítja.

\section{Frontend-megvalósítás}
\label{sec:ads-frontend}

\subsection{Listaoldal és szűrés}

Az \texttt{AdsListPage} a szűrőket (keresőmező, kategória-legördülő, ár-intervallum, „Saját
hirdetéseim” kapcsoló) komponens-lokális állapotban tartja, a találatokat pedig a React Query
kezeli. A lekérdezés kulcsa a teljes szűrőállapotot tartalmazza
(\texttt{["ads", categoryId, search, minPrice, maxPrice, showMyAds, page]}), így bármely szűrő
változása automatikusan új lekérdezést vált ki, a korábbi szűrőkombinációk eredménye pedig a
gyorsítótárból azonnal visszatölthető. Egy effektus gondoskodik róla, hogy szűrőváltáskor a lapozó
az első oldalra ugorjon. A „Saját hirdetéseim” kapcsoló a bejelentkezett felhasználó azonosítóját
adja át \texttt{ownerUserId} paraméterként, vagyis a szűrés ilyenkor is szerveroldalon történik. A
lapszámot a kliens a válasz \texttt{total} mezőjéből számítja, kilences lapmérettel.

\subsection{Hirdetésfeladás és az űrlap}

Az új hirdetés oldala (\texttt{AdCreatePage}) vendég látogatót a bejelentkezési oldalra irányít,
az eredeti célútvonalat útvonalállapotban megőrizve. Az űrlapot a létrehozó és a szerkesztő oldal
által közösen használt \texttt{AdForm} komponens adja: kontrollált beviteli mezőkkel dolgozik, a
kategórialistát a gyorsítótárazott \texttt{["ad-categories"]} lekérdezésből tölti, az új képeket
pedig memóriában gyűjti -- a fájlok háromféleképpen adhatók hozzá (fájlböngésző, húzd-és-ejtsd,
vágólap-beillesztés), és csak \texttt{image/} MIME-típusú fájlokat fogad el. Az előnézetekhez a
komponens objektum-URL-eket képez, amelyeket leszereléskor fel is szabadít. Szerkesztéskor a már
feltöltött képek nem azonnal törlődnek, hanem eltávolításra jelölhetők, és csak az űrlap
beküldésekor kerülnek ténylegesen törlésre.

A beküldés folyamata kétfázisú, mert a backend a hirdetést és a képeket külön végpontokon
kezeli: először a hirdetés jön létre, majd a kapott azonosítóval egyenként töltődnek fel a képek,
végül a kliens érvényteleníti a hirdetéslista gyorsítótárát és a részletező oldalra navigál
(\ref{lst:ads-create}.~kódrészlet).

\begin{lstlisting}[caption={A hirdetésfeladás kétfázisú beküldése az \texttt{AdCreatePage} komponensben (kivonat)}, label={lst:ads-create}]
const created = await mutation.mutateAsync(payload);

if (newFiles.length > 0) {
  setIsProcessingImages(true);
  for (const file of newFiles) {
    await adsApi.uploadImage(created.adId, file);
  }
}

await queryClient.invalidateQueries({ queryKey: ["ads"] });
await queryClient.invalidateQueries({ queryKey: ["ads", created.adId] });
navigate(`/ads/${created.adId}`, { replace: true });
\end{lstlisting}

\subsection{Részletező és szerkesztő oldal}

A részletező oldal (\texttt{AdDetailPage}) az útvonal-paraméterből kiolvasott azonosítóval, az
\texttt{["ads", adId]} kulcson kérdezi le a hirdetést. A kezelőgombok megjelenítéséről a
\texttt{canManage} feltétel dönt: a felhasználó adminisztrátor, vagy a hirdetés
\texttt{ownerUserId} mezője egyezik az azonosítójával -- ez ugyanaz a szabály, amelyet a backend
is kikényszerít, így a felület nem ígér olyan műveletet, amelyet a szerver elutasítana. A törlés
natív megerősítő kérdés után fut, sikere a lista-gyorsítótár érvénytelenítésével és a listaoldalra
navigálással zárul. Képtörléskor a kliens optimalizált frissítést alkalmaz: a részletező oldal
gyorsítótárazott adatából azonnal kiszűri a törölt képet, és csak ezután érvényteleníti a
listákat. A szerkesztő oldal (\texttt{AdEditPage}) ugyanazt az \texttt{AdForm} komponenst tölti
fel a meglévő hirdetés adataival; beküldéskor a módosítás mellett az újonnan csatolt képeket
feltölti, az eltávolításra jelölteket törli, majd szintén a részletező oldalra tér vissza.

% TODO: SC-04 képernyőkép (figures/sc-04-ad-form.png) — képelőnézettel és drag&drop zónával.
\begin{figure}[htbp]
	\centering
	\fbox{\parbox[c][5cm][c]{0.85\textwidth}{\centering SC-04: A hirdetésfeladó űrlap (elkészítendő)}}
	\caption{A hirdetésfeladó űrlap képelőnézettel és húzd-és-ejtsd feltöltési zónával}
	\label{fig:ad-form-screen}
\end{figure}

\section{Adatfolyam, biztonság és korlátok}
\label{sec:ads-adatfolyam}

A hirdetésfeladás teljes, képfeltöltéssel kombinált adatfolyamát
a~\ref{fig:ad-create-sequence}.~ábra szemlélteti. A folyamat jól példázza a REST-erőforrások
szétválasztásából adódó tervezési kompromisszumot: a hirdetés és a képek külön erőforrások, ezért
a kliens felelőssége a műveletsor levezénylése. Előnye, hogy a képfeltöltés hibája nem veszíti el
magát a hirdetést; hátránya, hogy a folyamat nem atomi -- ha a feltöltési ciklus félbeszakad, a
hirdetés a képek egy részével marad meg.

% TODO: F-01 szekvenciadiagram (figures/f-01-ad-create-sequence.pdf) — AdForm → POST /ads →
%       AdService → SaveChanges → képenként POST /ads/{id}/images → wwwroot/uploads.
\begin{figure}[htbp]
	\centering
	\fbox{\parbox[c][6.5cm][c]{0.9\textwidth}{\centering F-01: A hirdetésfeladás szekvenciadiagramja képfeltöltéssel (elkészítendő)}}
	\caption{A hirdetésfeladás és képfeltöltés szekvenciadiagramja}
	\label{fig:ad-create-sequence}
\end{figure}

Biztonsági szempontból a funkció többrétegű védelmet valósít meg: a módosító végpontok
tulajdonos- vagy adminisztrátor-ellenőrzéshez kötöttek; a tulajdonlás kettős invariánsát az
üzleti réteg és adatbázis-megszorítás együtt őrzi; a képfeltöltést kiterjesztés-fehérlista,
méretkorlát és generált fájlnév védi. A korlátok közé tartozik, hogy a kiterjesztés-ellenőrzés a
fájlnéven alapul, a fájl tényleges tartalmát (mágikus bájtok) a rendszer nem vizsgálja; a
feltöltött fájlok autentikáció nélkül, statikus tartalomként érhetők el, ami nyilvános piactérnél
elfogadható, de érzékeny tartalomnál nem volna az. Megjegyzendő továbbá, hogy a végpontok magát az
EF-entitást fogadják és adják vissza kérés- és válaszmodellként; ez kevesebb kódot igényel, mint a
dedikált DTO-réteg, cserébe a módosító kérések a szükségesnél több mezőt tartalmazhatnak -- a
kontroller ezért másolja át mezőnként, explicit listával a módosítható értékeket. Végül a
kliensoldali állapotlista jelenleg rögzített \texttt{ACTIVE} alapértelmezettel dolgozik, és bár az
állapotlekérdező végpont és kliensfüggvény létezik, a vizsgált oldalak nem használják -- az
állapot alapú szűrés így a felületen kiaknázatlan lehetőség maradt.

% !TeX root = ../main.tex
% 7. fejezet — Vadásztársaságok és tagságkezelés
% Hivatkozott források: Backend/HUNting.API/Controllers/SocietiesController.cs, NewsController.cs,
%   Services/SocietyService.cs, Models/Societies/*, Data/HuntingContext.cs, Data/SeedData.cs,
%   Frontend/HUNting/src/services/societiesApi.ts, pages/Societies/SocietiesPage.tsx,
%   SocietyDetailPage.tsx, SocietyManagementPage.tsx, MySocietiesPage.tsx

\chapter{Vadásztársaságok és tagságkezelés}
\label{chap:tarsasagok}

A vadásztársasági funkciócsoport a HUNting közösségi magja: itt találkozik a két fióktípus. Az
egyéni vadászok társaságokat kereshetnek és tagfelvételüket kérhetik, a társasági fiókok pedig a
kérelmek elbírálásával, a taglistával és a belső bejegyzésekkel igazgatják közösségüket. A fejezet
a tagfelvételi kérelem életciklusára mint a rendszer legjellegzetesebb üzleti folyamatára
összpontosít, zárásként pedig tisztázza a hírfolyam (NewsFeed) jelenlegi, részlegesen megvalósított
állapotát.

\section{A funkció célja és felhasználói működése}
\label{sec:soc-cel}

A társaságkatalógus nyilvános: a látogató név szerint szűrhető listában böngészheti a regisztrált
társaságokat, és megnyithatja bármelyik profilját, amely a bemutatkozást, az elérhetőségeket, az
alapítás dátumát és a taglétszámot mutatja. Bejelentkezett, egyéni fiókkal rendelkező felhasználó a
profiloldalról tagfelvételi kérelmet nyújthat be, opcionális bemutatkozó üzenettel. A kérelemről a
társaság tulajdonosa dönt a társaságkezelő felületen, ahol a függőben lévő jelentkezések, a
taglista és a belső posztok is egy helyen érhetők el. Elfogadás után az új tag a „Társaságaim”
oldalon látja tagságait, és társaságot választva olvashatja annak belső bejegyzéseit -- a posztok
tehát csak a tagság megszerzése után válnak láthatóvá.

% TODO: SC-05 és SC-06 képernyőkép (figures/sc-05-societies.png, sc-06-society-detail.png)
\begin{figure}[htbp]
	\centering
	\fbox{\parbox[c][5cm][c]{0.85\textwidth}{\centering SC-05/SC-06: A társaságkatalógus és a társasági profil a jelentkezési űrlappal (elkészítendő)}}
	\caption{A társaságkatalógus és egy társaság profilja a tagfelvételi űrlappal}
	\label{fig:societies-screens}
\end{figure}

\section{Adatmodell}
\label{sec:soc-adatmodell}

A funkciócsoport öt entitásra épül (\texttt{Models/Societies/}). A \texttt{Society} a társaság
törzsadatait tárolja (egyedi indexszel védett e-mail-cím, név, telefonszám, bemutatkozás,
létrehozási idő), és a~\ref{chap:felhasznalok}.~fejezetben bemutatott módon egy-egy kapcsolatban
áll tulajdonos felhasználójával. A társaságon belüli rangokat a \texttt{SocietyRole} szótártábla
írja le; a kezdeti adatfeltöltés négy szerepet hoz létre (Vezeto, Gazda, Vador, Tag).

A tagságot a \texttt{SocietyMember} kapcsolótábla valósítja meg, amely a klasszikus több-a-többhöz
mintát követi, de a puszta kapcsolaton túl attribútumokat is hordoz: a tag társaságon belüli
szerepét és a csatlakozás időpontját. Elsődleges kulcsa összetett -- a
\texttt{(SocietyId, UserId)} pár --, ami adatbázisszinten zárja ki, hogy ugyanaz a felhasználó
kétszer legyen tagja ugyanannak a társaságnak (\ref{lst:member-config}.~kódrészlet). Ugyanilyen
összetett kulcsú a \texttt{SocietyMembershipRequest} entitás is, vagyis egy felhasználóhoz
társaságonként legfeljebb egy kérelemrekord tartozhat; az ebből következő „újranyitási” viselkedést
a~\ref{sec:soc-backend}.~alfejezet mutatja be. A kérelem állapotát a
\texttt{MembershipRequestStatus} felsorolt típus írja le (\texttt{Pending}, \texttt{Approved},
\texttt{Rejected}), emellett a rekord az elbírálás időpontját és az elbíráló azonosítóját is
rögzíti.

\begin{lstlisting}[language={[Sharp]C}, caption={A tagsági kapcsolótábla konfigurációja a \texttt{HuntingContext} osztályban (kivonat)}, label={lst:member-config}]
modelBuilder.Entity<SocietyMember>()
    .HasKey(x => new { x.SocietyId, x.UserId });

modelBuilder.Entity<SocietyMember>()
    .HasOne(x => x.Society)
    .WithMany(x => x.Members)
    .HasForeignKey(x => x.SocietyId)
    .OnDelete(DeleteBehavior.Cascade);

modelBuilder.Entity<SocietyMember>()
    .HasOne(x => x.User)
    .WithMany(x => x.Societies)
    .HasForeignKey(x => x.UserId)
    .OnDelete(DeleteBehavior.Cascade);

modelBuilder.Entity<SocietyMember>()
    .HasOne(x => x.Role)
    .WithMany(x => x.Members)
    .HasForeignKey(x => x.RoleId)
    .OnDelete(DeleteBehavior.Restrict);
\end{lstlisting}

A törlési szabályok a rekord jellegéhez igazodnak: a tagság a társasággal és a felhasználóval
együtt kaszkádolva törlődik, a szerephivatkozás viszont \texttt{Restrict}, így használatban lévő
társasági szerep nem távolítható el. A kérelemnél árnyaltabb a kép: a társaság törlése a
kérelmeket is viszi, a kérelmező felhasználó törlése tiltott, amíg kérelme létezik
(\texttt{Restrict}), az elbíráló törlésekor pedig a hivatkozás \texttt{SetNull} értékre vált --
a döntés ténye tehát az elbíráló fiókjának megszűnése után is megmarad. A belső bejegyzéseket a
\texttt{SocietyPost} entitás tárolja (cím, tartalom, opcionális képútvonal, szerző, időbélyeg),
amely a társasággal és a szerzővel együtt kaszkádolva törlődik.

% TODO: E-02 osztálydiagram (figures/e-02-society-domain.pdf) — ApplicationUser, Society,
%       SocietyMember, SocietyMembershipRequest, SocietyRole, SocietyPost kapcsolatokkal.
\begin{figure}[htbp]
	\centering
	\fbox{\parbox[c][6.5cm][c]{0.9\textwidth}{\centering E-02: A társasági domain UML-osztálydiagramja (elkészítendő)}}
	\caption{A társasági domain entitásai és kapcsolataik}
	\label{fig:society-domain}
\end{figure}

\section{Backend-megvalósítás}
\label{sec:soc-backend}

A funkciót a \texttt{SocietiesController} és a \texttt{SocietyService} párosa valósítja meg. A
szolgáltatás minden adatelérési műveletét közvetlenül a \texttt{HuntingContext}-en végzi: a
társaság-törzsadatok egyszerű CRUD-műveleteit éppúgy, mint az összetettebb lekérdezéseket és a
kapcsolótáblák kezelését. A végpontokat
a~\ref{tab:soc-endpoints}.~táblázat összegzi.

\begin{table}[htbp]
	\centering
	\caption{A \texttt{SocietiesController} végpontjai}
	\label{tab:soc-endpoints}
	\begin{tabular}{|l|l|p{5.2cm}|}
		\hline
		\textbf{Végpont} & \textbf{Hozzáférés} & \textbf{Feladat} \\
		\hline
		GET \texttt{/api/Societies} & anonim & Társaságok listája \\
		\hline
		GET \texttt{/api/Societies/\{id\}} & anonim & Egy társaság részletei \\
		\hline
		GET \texttt{/api/Societies/my} & bejelentkezett & A hívó tagságai \\
		\hline
		GET \texttt{/api/Societies/\{id\}/members} & tulajdonos v. admin & Taglista \\
		\hline
		POST \texttt{/api/Societies/\{id\}/join} & bejelentkezett Person & Tagfelvételi kérelem \\
		\hline
		GET \texttt{/api/Societies/\{id\}/requests} & tulajdonos v. admin & Függő kérelmek \\
		\hline
		POST \texttt{.../requests/\{userId\}/approve} & tulajdonos v. admin & Kérelem elfogadása \\
		\hline
		POST \texttt{.../requests/\{userId\}/reject} & tulajdonos v. admin & Kérelem elutasítása \\
		\hline
		GET \texttt{/api/Societies/\{id\}/posts} & tag, tulajdonos v. admin & Belső posztok \\
		\hline
		POST \texttt{/api/Societies/\{id\}/posts} & tulajdonos v. admin & Poszt létrehozása képpel \\
		\hline
		DELETE \texttt{.../posts/\{postId\}} & tulajdonos v. admin & Poszt törlése \\
		\hline
	\end{tabular}
\end{table}

\subsection{A tagfelvételi kérelem életciklusa}

A jelentkezési végpont először a fióktípust ellenőrzi: a token \texttt{account\_type} claimjéből
megállapítja, hogy a hívó társasági fiók-e, és ha igen, elutasítja a kérést -- társaság nem
csatlakozhat másik társasághoz. Ezután a szolgáltatás a kérelem állapotgépét hajtja végre
(\ref{lst:join-lifecycle}.~kódrészlet): ha a hívó már tag, a kérés hibával zárul; ha van függőben
lévő kérelme, szintén; ha viszont van korábbi -- elfogadott vagy elutasított -- kérelemrekordja,
azt a rendszer \emph{újranyitja}: az állapotot \texttt{Pending} értékre állítja, frissíti az
üzenetet és a beadás időpontját, az elbírálási mezőket pedig kiüríti. Új rekord csak akkor jön
létre, ha a felhasználónak még egyáltalán nem volt kérelme az adott társaságnál -- ezt a
megoldást az összetett kulcs kényszeríti ki, amely társaságonként és felhasználónként egyetlen
kérelemrekordot enged.

\begin{lstlisting}[language={[Sharp]C}, caption={A kérelem létrehozása és újranyitása a \texttt{SocietyService} osztályban (kivonat)}, label={lst:join-lifecycle}]
var existingMember = await _ctx.SocietyMembers
    .AnyAsync(m => m.SocietyId == societyId && m.UserId == userId);
if (existingMember)
{
    throw new InvalidOperationException("Mar tagja vagy a tarsasagnak.");
}

var request = await _ctx.SocietyMembershipRequests.FindAsync(societyId, userId);
if (request != null)
{
    if (request.Status == MembershipRequestStatus.Pending)
    {
        throw new InvalidOperationException("Mar van fuggoben levo kerelmed.");
    }

    request.Status = MembershipRequestStatus.Pending;
    request.Message = message;
    request.CreatedAt = DateTime.UtcNow;
    request.ProcessedAt = null;
    request.ProcessedByUserId = null;
}
else
{
    request = new SocietyMembershipRequest { /* ... */ };
    _ctx.SocietyMembershipRequests.Add(request);
}
\end{lstlisting}

Az elbírálás a tulajdonos vagy adminisztrátor kiváltsága, amit a kontroller közös
\texttt{IsOwnerOrAdmin} segédmetódusa ellenőriz. Az elfogadó művelet idempotens: már elfogadott
kérelemre nem fut le újra. Sikeres elfogadáskor a kérelem állapota \texttt{Approved} lesz, a
rekord rögzíti az elbírálót és az időpontot, majd -- ha a felhasználó időközben nem lett tag --
létrejön a tagsági rekord az alapértelmezett „Tag” társasági szereppel (ennek hiányában az első
elérhető szereppel). Az elutasítás csak az állapotot és az elbírálási mezőket írja; a korábban
bemutatott újranyitási szabály miatt az elutasított jelentkező később újra próbálkozhat. A teljes
életciklust a~\ref{fig:request-lifecycle}.~ábra foglalja össze.

% TODO: F-02 állapotdiagram (figures/f-02-request-lifecycle.pdf) — (nincs) → Pending →
%       Approved / Rejected → újranyitás Pendingre.
\begin{figure}[htbp]
	\centering
	\fbox{\parbox[c][5.5cm][c]{0.9\textwidth}{\centering F-02: A tagfelvételi kérelem állapotdiagramja (elkészítendő)}}
	\caption{A tagfelvételi kérelem életciklusa}
	\label{fig:request-lifecycle}
\end{figure}

\subsection{Belső bejegyzések}

A posztok olvasási joga szélesebb, mint az írási: a listát a társaság bármely tagja, a tulajdonos
és az adminisztrátor kérheti le -- a kontroller a tagságot és a tulajdonlást közvetlen
adatbázis-lekérdezéssel ellenőrzi --, új posztot viszont csak a tulajdonos vagy adminisztrátor
hozhat létre. A poszt többrészes kérésben érkezik (cím, tartalom, opcionális kép, 10~MiB
kéréslimittel), a képkezelés pedig a hirdetéseknél megismert mintát követi:
kiterjesztés-fehérlista, GUID fájlnév, \texttt{uploads/societies} célkönyvtár. Poszt törlésekor a
kontroller ellenőrzi, hogy a bejegyzés valóban az útvonalban megadott társasághoz tartozik-e,
majd a képfájlt és a rekordot is eltávolítja.

\section{Frontend-megvalósítás}
\label{sec:soc-frontend}

\subsection{Katalógus és jelentkezés}

A katalógusoldal (\texttt{SocietiesPage}) a teljes társaságlistát kéri le, és a névszűrést --
a lista mérete miatt elfogadható kompromisszumként -- kliensoldalon végzi. A profiloldal
(\texttt{SocietyDetailPage}) a \texttt{["society", societyId]} kulcson kérdezi le az adatokat, és
a jelentkezési űrlapot csak akkor jeleníti meg, ha a látogató bejelentkezett és fióktípusa
\texttt{Person} -- ugyanaz a szabály, amelyet a backend a token claimjéből kikényszerít. A
jelentkezés mutációja sikeres beküldés után visszaigazoló üzenetet jelenít meg és üríti az
üzenetmezőt; a szerver által dobott üzleti hibák (már tag, már van függő kérelem) az űrlap
hibaszövegeként jelennek meg.

\subsection{Társaságkezelő és tagsági nézet}

A társaságkezelő oldal (\texttt{SocietyManagementPage}) hozzáférési feltétele kliensoldalon is
kettős: bejelentkezett, \texttt{Society} típusú fiók kell, amelyhez tartozik birtokolt társaság
(\texttt{user.ownedSociety}) -- ennek hiányában az oldal magyarázó üzenetet jelenít meg. A
társaság azonosítóját a tárolt profil \texttt{ownedSociety} mezőjéből veszi, és három párhuzamos
lekérdezést indít: tagok (\texttt{["society-members", id]}), függő kérelmek
(\texttt{["society-requests", id]}) és posztok (\texttt{["society-posts", id]}). A mutációk
célzottan érvénytelenítik az érintett gyorsítótárakat: az elfogadás a kérelmek \emph{és} a tagok
listáját is frissíti -- hiszen mindkettőt megváltoztatja --, az elutasítás csak a kérelmekét, a
posztműveletek pedig a posztokét. Az új poszt űrlapja cím- és tartalommezőből, valamint natív
fájlválasztóból áll; beküldés után a mezők kiürülnek.

A tagsági nézet (\texttt{MySocietiesPage}) az egyéni fiókok tükörképe: a \texttt{["my-societies"]}
lekérdezés a hívó tagságait adja vissza, a kiválasztott társaság posztjait pedig egy feltételesen
engedélyezett lekérdezés (\texttt{enabled: Boolean(selectedSociety)}) tölti be -- így hálózati
kérés csak akkor indul, amikor a felhasználó ténylegesen társaságot választott. A posztok képei a
\texttt{resolveAssetUrl} segédfüggvényen keresztül kapják meg a backend abszolút címét.

% TODO: SC-07 és SC-08 képernyőkép (figures/sc-07-society-manage.png, sc-08-my-societies.png)
\begin{figure}[htbp]
	\centering
	\fbox{\parbox[c][5cm][c]{0.85\textwidth}{\centering SC-07/SC-08: A társaságkezelő és a „Társaságaim” oldal (elkészítendő)}}
	\caption{A társaságkezelő felület a függő kérelmekkel, valamint a tagsági nézet a belső posztokkal}
	\label{fig:society-manage-screens}
\end{figure}

\section{A hírfolyam (NewsFeed) állapota és a funkciócsoport korlátai}
\label{sec:soc-news}

A kódbázis a belső posztok mellett egy második, hasonló rendeltetésű entitáscsaládot is tartalmaz:
a \texttt{NewsFeed} hírbejegyzést és a hozzá tartozó \texttt{NewsVisibility} láthatósági
szótártáblát (\texttt{MEMBERS}, \texttt{PUBLIC}, \texttt{OFFICERS} kódokkal), két mintahírrel a
kezdeti adatfeltöltésben. Fontos azonban őszintén rögzíteni a megvalósultság tényleges szintjét:
a \texttt{NewsController} teljes CRUD-ot kínál -- az író végpontjait (létrehozás, módosítás,
törlés) az \texttt{[Authorize]} attribútum adminisztrátori szerep-megkötése védi, az olvasó
végpontok anonim hívhatók --, a tárolt láthatósági kód szerinti
szűrés azonban nem valósul meg, a frontenden pedig a hírekhez semmilyen oldal vagy API-hívás nem
tartozik. A hírfolyam tehát \emph{backend-oldali előkészítésnek}
tekintendő, nem működő felhasználói funkciónak; a dolgozat a tagok felé irányuló tényleges
kommunikációs csatornaként a \texttt{SocietyPost} bejegyzéseket dokumentálja. A hiányzó
láthatóságszűrést és felhasználói felületet a továbbfejlesztési irányok között
(\ref{chap:osszefoglalas}.~fejezet) kell rendezni.

A funkciócsoport további ismert korlátai: a katalógusoldal kiemelt „Jelentkezés tagságra” gombjához
jelenleg nem tartozik eseménykezelő -- a működő jelentkezési űrlap a profiloldalon él --; a
taglistáról nincs tageltávolító felületi művelet (a szolgáltatásban a metódus elkészült, végpont
nem tartozik hozzá); a társasági szerepek (Gazda, Vador) pedig a tagrekordban tárolódnak ugyan, de
jogosultsági következményük jelenleg nincs -- a döntéseket minden ponton a tulajdonos-vagy-admin
szabály vezérli.

% !TeX root = ../main.tex
% 8. fejezet — Vadásznapló
% Hivatkozott források: Backend/HUNting.API/Controllers/HuntingLogsController.cs,
%   Models/Hunting/HuntingLog.cs, Data/HuntingContext.cs,
%   Frontend/HUNting/src/services/huntingLogsApi.ts, pages/HuntingLogs/HuntingLogsPage.tsx

\chapter{Vadásznapló}
\label{chap:naplo}

A vadásznapló a hagyományos, papíralapú lőjegyzék digitális megfelelője: az egyéni vadász saját,
személyes nyilvántartása a vadászatairól. A funkció a gerinc legkompaktabb tagja, ugyanakkor jól
példázza a fióktípus szerinti korlátozást és az adminisztrátori hatókör-kiterjesztést, ezért a
dolgozat külön fejezetben tárgyalja.

\section{A funkció célja és felhasználói működése}
\label{sec:log-cel}

A napló kizárólag bejelentkezett, egyéni (\texttt{Person}) fiókkal érhető el: vendég látogatót az
oldal a bejelentkezésre irányít, társasági fióknak pedig magyarázó üzenetet jelenít meg. A
felhasználó időrendben -- a vadászat időpontja szerint csökkenő sorrendben -- látja saját
bejegyzéseit, amelyek kártyái lenyitva mutatják a részleteket: helyszín, vadfaj, jegyzetek és a
csatolt fénykép. Új bejegyzés a lista fölötti űrlapon rögzíthető: kötelező a vadászat időpontja,
opcionális a helyszín, a vadfaj, a jegyzet és egy fénykép, amelyhez a felület azonnali előnézetet
ad. A bejegyzések szigorúan személyesek: minden felhasználó csak a sajátjait látja, az
adminisztrátor azonban -- kizárólag az API \texttt{scope=all} paraméterével -- a teljes
naplóállományt is lekérdezheti.

% TODO: SC-10 képernyőkép (figures/sc-10-hunting-logs.png) — lista + nyitott űrlap képelőnézettel.
\begin{figure}[htbp]
	\centering
	\fbox{\parbox[c][5cm][c]{0.85\textwidth}{\centering SC-10: A vadásznapló oldala (elkészítendő)}}
	\caption{A vadásznapló bejegyzéslistája és az új bejegyzés űrlapja}
	\label{fig:hunting-logs-screen}
\end{figure}

\section{Adatmodell}
\label{sec:log-adatmodell}

A funkciót egyetlen entitás szolgálja ki, a \texttt{HuntingLog}
(\texttt{Models/Hunting/HuntingLog.cs}), amely az adatbázisban \texttt{LogId} oszlopnévre képezett
kulccsal rendelkezik. A bejegyzés a vadászat időpontját (\texttt{OccurredAt}), a helyszín szöveges
megnevezését és opcionális koordinátáit, a vadfajt (\texttt{AnimalType}), a lövés tényét
(\texttt{ShotFired}), szabad szöveges jegyzetet, a csatolt kép útvonalát, valamint a társasági
megosztás szándékát jelző \texttt{ShareWithSociety} kapcsolót tárolja, a rögzítés időbélyegével
együtt. A bejegyzés a felhasználóhoz kötelező idegen kulccsal kapcsolódik, és a fiók törlésekor
kaszkádolva törlődik -- a napló a fiók része, önálló életciklusa nincs. Más entitással a
naplónak nincs kapcsolata, ami a funkció személyes jellegét az adatmodell szintjén is tükrözi.

\section{Backend-megvalósítás}
\label{sec:log-backend}

A \texttt{HuntingLogsController} a kisebb kontrollerek közé tartozik, és -- az
\texttt{AuthController}-hez hasonlóan -- nem dedikált szolgáltatáson, hanem közvetlenül a
\texttt{HuntingContext}-en dolgozik; a rétegzés e pontján a projekt az egyszerűséget választotta a
szigorú egységesség helyett. Mind az öt végpontja \texttt{[Authorize]} attribútummal védett.

A listázó végpont alapértelmezetten a hívó saját bejegyzéseire szűr. A hatókör kiterjesztése két
feltétel együttes teljesüléséhez kötött: a hívónak \emph{Administrator} szerepben kell lennie,
\emph{és} a kérésben kis- és nagybetűtől függetlenül meg kell adnia a \texttt{scope=all}
paramétert -- adminisztrátor paraméter nélkül tehát szintén csak a saját naplóját látja. Az
egyedi lekérdezés, a módosítás és a törlés a már ismert erőforrás-tulajdonos mintát követi:
tulajdonos vagy adminisztrátor férhet hozzá, egyébként a válasz 403-as.

A létrehozó végpont többrészes űrlapot fogad (\texttt{HuntingLogCreateRequest}, 10~MiB
kéréslimittel), és két ponton tér el a sablontól. Egyrészt fióktípus-korlátot érvényesít: a token
\texttt{account\_type} claimje alapján a társasági fiók kérését hibával utasítja el, mivel a
naplóvezetés az egyéni vadász tevékenysége. Másrészt a kép mentése a bejegyzés mentése
\emph{előtt} történik, így érvénytelen fájl esetén az adatbázis érintetlen marad; a képkezelés a
közös mintát követi (kiterjesztés-fehérlista, GUID fájlnév, \texttt{uploads/logs} célkönyvtár).
Mentés után a kontroller betölti a bejegyzés felhasználó-navigációját, hogy a 201-es válasz a
megjelenítéshez szükséges szerzőadatokat is tartalmazza. Törléskor -- a hirdetésekhez hasonlóan --
a lemezen tárolt képfájl is törlődik, így nem marad árva állomány a webgyökérben.

\section{Frontend-megvalósítás, adatfolyam és korlátok}
\label{sec:log-frontend}

A \texttt{HuntingLogsPage} a hozzáférést két lépcsőben ellenőrzi: hitelesítetlen látogatót
átirányít, nem \texttt{Person} fiókot szöveges üzenettel zár ki, és a \texttt{["huntingLogs"]}
lekérdezést is csak jogosult felhasználónál engedélyezi (\texttt{enabled} feltétel). Az űrlap
kontrollált mezőkkel dolgozik; az időpontmező \texttt{datetime-local} típusú, kezdőértéke a
helyi időzónára igazított aktuális idő, beküldéskor pedig ISO-formátumra alakul. A fájlválasztó
csak \texttt{image/} MIME-típusú fájlt fogad el, az előnézethez objektum-URL-t képez, és azt a
csere vagy a komponens leszerelése után felszabadítja.

A beküldés a \texttt{huntingLogsApi.create} függvényen át \texttt{FormData} tartalomként utazik,
a mezőneveket a backend kötési neveire képezve. Siker esetén a kliens nem érvénytelenítéssel,
hanem közvetlen gyorsítótár-frissítéssel reagál: az új bejegyzést a \texttt{setQueryData}
hívással a lista elejére fűzi, így a felület újabb hálózati kérés nélkül, azonnal mutatja a
friss adatot -- kis adatmennyiségnél ez gyorsabb és olcsóbb, mint az újratöltés. Ezt követően az
űrlap bezárul és alaphelyzetbe áll.

A funkció ismert korlátai jól körülhatárolhatók. A kliens a \texttt{ShotFired} és a
\texttt{ShareWithSociety} mezőt jelenleg rögzített \texttt{false} értékkel küldi -- a felületen
nincs hozzájuk vezérlőelem --, és a \texttt{ShareWithSociety} jelzőt a rendszer más pontja sem
dolgozza fel: a társasági megosztás az adatmodellben előkészített, de végponttal és felülettel még
nem támogatott képesség. A backend módosító és törlő végpontjai készen állnak, ám a felület nem
kínál hozzájuk szerkesztő- vagy törlőműveletet, ahogy az adminisztrátori \texttt{scope=all}
lekérdezéshez sem tartozik nézet. Végül a listázás lapozás nélküli: a személyes napló várható
mérete mellett ez elfogadható, de hosszú távon a hirdetéseknél bemutatott lapozott minta
átvétele indokolt. Ezek a pontok a továbbfejlesztési irányok között
(\ref{chap:osszefoglalas}.~fejezet) jelennek meg újra.

% !TeX root = ../main.tex
% 9. fejezet — Térképes funkciók
% Hivatkozott források: Backend/HUNting.API/Controllers/MapController.cs, Services/MapService.cs,
%   Models/Map/*, Data/HuntingContext.cs, Frontend/HUNting/src/services/mapApi.ts,
%   pages/Map/MapPage.tsx

\chapter{Térképes funkciók}
\label{chap:terkep}

A térképmodul a HUNting térinformatikai komponense: megfigyelőpontok, etetők, veszélyhelyek és
találkozási pontok, valamint vadászterület-határok rögzítését teszi lehetővé személyes és
társasági rétegeken. A fejezet a kliensoldali Google Maps-integráció és a szerveroldali
hatókör-kezelés (scope) együttműködését mutatja be -- ez a rendszer legösszetettebb
kliens--szerver együttműködése.

\section{A funkció célja és felhasználói működése}
\label{sec:map-cel}

A térképoldal Magyarország-középpontú Google-térképpel indul. A bejelentkezett felhasználó
rétegválasztóból dönt: a „Saját jelölések” réteg a kizárólag általa rögzített, személyes
objektumokat mutatja, emellett minden olyan társaság közös rétege is választható, amelynek tagja.
Marker felvételéhez a felhasználó jelölő módba vált, a térképre kattintva kijelöli a pontot, majd
ikontípust (például etető, megfigyelőpont), jegyzetet és láthatósági kódot rendel hozzá. Terület
rajzolásához a rajzoló mód a Google rajzeszközét aktiválja: a csúcspontok lehelyezése után a
terület nevet, leírást, körvonal- és kitöltőszínt, valamint láthatóságot kaphat. A rögzített
objektumok a térkép melletti listában is megjelennek, onnan fókuszba hozhatók, és -- saját
objektum esetén -- törölhetők; az információs ablak a jegyzetet, illetve a terület leírását
mutatja.

% TODO: SC-09 képernyőkép (figures/sc-09-map.png) — markerek + polygon-rajzolás közben; Google Maps kulcs kell.
\begin{figure}[htbp]
	\centering
	\fbox{\parbox[c][5.5cm][c]{0.85\textwidth}{\centering SC-09: A térképoldal markerekkel és területrajzolással (elkészítendő)}}
	\caption{A térképoldal: személyes réteg markerekkel és rajzolás alatt álló területtel}
	\label{fig:map-screen}
\end{figure}

\section{Adatmodell}
\label{sec:map-adatmodell}

A térképi objektumokat hat entitás írja le (\texttt{Models/Map/}). A \texttt{MapMarker} egy
pontszerű jelölés: koordinátapárját a rendszer \texttt{decimal} típusban tárolja (a lebegőpontos
kerekítési hibák elkerülésére), ezen felül ikontípus-hivatkozást, opcionális jegyzetet,
láthatósági kódot, valamint létrehozó- és opcionális társasági hivatkozást hordoz. Az
ikontípusokat a \texttt{MapIconType} szótártábla adja (név, alapértelmezett szín, glif); a seed
négy típust rögzít (Megfigyelopont, Eteto, Veszely, Talalkozo). A láthatósági kódok a
\texttt{MapVisibility} referenciatáblában élnek (\texttt{PUBLIC}, \texttt{MEMBERS},
\texttt{ADMIN}), karakterláncos kulccsal.

A területeket a \texttt{MapPolygon} entitás és a hozzá tartozó \texttt{MapPolygonPoint}
pontrekordok együttese írja le: a sokszög neve, leírása, színei és láthatósága mellett a
csúcspontok külön táblában, \texttt{Seq} sorszámmal tárolódnak. A
\texttt{(PolygonId, Seq)} párra egyedi index vonatkozik, így a csúcspontok sorrendje -- amely a
sokszög alakját meghatározza -- adatbázisszinten is egyértelmű; a pontok a sokszöggel együtt
kaszkádolva törlődnek. A létrehozó felhasználóra mutató idegen kulcs minden objektumnál
\texttt{Restrict} viselkedésű, a társasági hivatkozás viszont kaszkád: a társaság törlése a közös
rétegét is magával viszi.

A kódbázis tartalmaz egy harmadik objektumcsaládot is: a \texttt{MapPath} útvonal-entitást és
\texttt{MapPathPoint} pontjait, a sokszögekével azonos felépítésben, kész migrációval. Ezekhez
azonban \emph{sem végpont, sem kliensoldali funkció nem tartozik} -- az útvonalrögzítés tehát az
adatmodell szintjén előkészített, de meg nem valósított képesség, amelyet a dolgozat a
továbbfejlesztési irányok között tart számon.

\section{Backend-megvalósítás}
\label{sec:map-backend}

A \texttt{MapController} minden adatműveletet az \texttt{IMapService} szolgáltatásra bíz, és --
e projektben egyedüliként -- a lekérdezésekhez \texttt{CancellationToken}-t is továbbít, így a
kliens által megszakított kérések a szerveren sem futnak feleslegesen tovább. A
\texttt{/api/Map/config} végpont a felület felépítéséhez szükséges referenciaadatokat (ikontípusok
és láthatósági kódok) adja vissza egyetlen hívásban.

\subsection{Hatókör-alapú olvasás}

Az olvasó végpontokon nincs \texttt{[Authorize]} attribútum: a hozzáférést a kérés
\texttt{scope} paramétere alapján, kézzel ellenőrzi a kontroller. A nyers paramétert egy kis
értéktípus normalizálja (\ref{lst:map-scope}.~kódrészlet), ismeretlen vagy hiányzó hatókörnél
pedig a válasz üres lista -- a végpont tehát anonim hívásra sem hibázik, de adatot sem ad ki.

\begin{lstlisting}[language={[Sharp]C}, caption={A hatókör-paraméter normalizálása (\texttt{MapLayerScope})}, label={lst:map-scope}]
public readonly record struct MapLayerScope(string Value)
{
    public static MapLayerScope None => new("none");
    public static MapLayerScope Personal => new("personal");
    public static MapLayerScope Society => new("society");

    public static MapLayerScope From(string? raw) =>
        raw?.ToLowerInvariant() switch
        {
            "personal" => Personal,
            "society" => Society,
            _ => None
        };
}
\end{lstlisting}

A \texttt{personal} hatókör bejelentkezést követel, és a szolgáltatás azokra az objektumokra
szűr, amelyeket a hívó hozott létre \emph{és} nem tartoznak társasághoz. A \texttt{society}
hatókörnél a kontroller a társaság-azonosító megléte után a tagságot is ellenőrzi
(\texttt{UserBelongsToSocietyAsync}); nem tag hívó 403-as választ kap. A társasági réteg tehát a
tagok közös nézete: a lekérdezés a társaság \emph{összes} objektumát visszaadja, függetlenül
attól, melyik tag rögzítette.

\subsection{Írási műveletek és a pontok cseréje}

Objektum létrehozásakor a kontroller a társasági réteghez kötést szintén tagsághoz köti, a
láthatósági kód alapértelmezése \texttt{PUBLIC}. A sokszög csúcspontjait a kérésben érkező
sorrend szerint, folytatólagos \texttt{Seq} sorszámmal képezi le pontrekordokká. Módosításkor és
törléskor a jogosultságot a szolgáltatás \texttt{CanModify} szabálya dönti el: kizárólag az
objektum \emph{létrehozója} módosíthat -- ez szigorúbb a többi funkciónál megismert
tulajdonos-vagy-admin mintánál, és azt jelenti, hogy a társasági rétegen minden tag csak a saját
jelöléseit kezelheti. Jogosulatlan kísérletnél a szolgáltatás \texttt{UnauthorizedAccessException}
kivételt dob, amelyet a kontroller 403-as válaszra fordít. A sokszögmódosítás a pontok teljes
cseréjével dolgozik (\ref{lst:map-points}.~kódrészlet): a meglévő pontrekordok törlése után az új
pontlista kerül be -- az egyszerű és mindig konzisztens megoldás, amely kis pontszámnál
hatékonyabb, mint a különbségalapú frissítés.

\begin{lstlisting}[language={[Sharp]C}, caption={A sokszög csúcspontjainak cseréje a \texttt{MapService.UpdatePolygonAsync} metódusban (kivonat)}, label={lst:map-points}]
existing.Name = polygon.Name;
existing.Description = polygon.Description;
existing.VisibilityCode = polygon.VisibilityCode;
existing.StrokeColor = polygon.StrokeColor;
existing.FillColor = polygon.FillColor;
existing.SocietyId = polygon.SocietyId;

_context.MapPolygonPoints.RemoveRange(existing.Points);
existing.Points = polygon.Points
    .Select(pp => new MapPolygonPoint
    {
        Seq = pp.Seq,
        Latitude = pp.Latitude,
        Longitude = pp.Longitude
    })
    .ToList();

await _context.SaveChangesAsync(cancellationToken);
\end{lstlisting}

A válaszok itt nem EF-entitások, hanem dedikált rekord-DTO-k (\texttt{MapMarkerResponse},
\texttt{MapPolygonResponse}), statikus \texttt{FromEntity} leképezőkkel; a sokszögválasz a
pontokat \texttt{Seq} szerint rendezve adja vissza -- a térképmodul e tekintetben a projekt
legletisztultabb API-felületét nyújtja.

\section{Frontend-megvalósítás}
\label{sec:map-frontend}

A \texttt{MapPage} a \texttt{@react-google-maps/api} könyvtárra épül: a Google Maps
JavaScript-API-t a \texttt{useLoadScript} horog tölti be a \texttt{drawing} kiegészítő
könyvtárral, magyar nyelvi és területi beállítással; az API-kulcsot a
\texttt{VITE\_GOOGLE\_MAPS\_API\_KEY} környezeti változó adja. A deklaratív komponenskészlet
(\texttt{MarkerF}, \texttt{PolygonF}, \texttt{InfoWindowF}, \texttt{DrawingManagerF}) a React
állapotából rajzolja a térképet.

A rétegválasztó bejegyzései a „Saját jelölések” rétegből és a felhasználó tagságaiból épülnek fel
(ez utóbbiakat a \texttt{societiesApi.my} hívás adja), a kiválasztás pedig a markerek és a
sokszögek lekérdezését vezérli: a React Query kulcsa a hatókört és a társaság-azonosítót is
tartalmazza (\texttt{["map", "markers", scope, societyId]}), a lekérdezés pedig csak akkor
engedélyezett, ha van API-kulcs és tényleges rétegválasztás. Rétegváltáskor az oldal minden
szerkesztési vázlatot és kijelölést visszaállít, így a módok nem keveredhetnek rétegek között.

A szerkesztés háromállapotú módvezérléssel működik (megtekintés, jelölő, rajzoló). Jelölő módban
a térképre kattintás koordinátáiból vázlat (draft) készül, amelyet a felhasználó ikontípussal,
jegyzettel és láthatósággal egészít ki; a mentés mutációja sikerkor a markerlekérdezéseket
érvényteleníti. Rajzoló módban a \texttt{DrawingManagerF} adja a csúcspontokat; a kliens csak
legalább három érvényes pontból álló sokszöget enged menteni. A megjelenítés előtt a kliens
védekezően normalizál: a koordinátákat számmá alakítja (a vesszős tizedesjelölést is kezelve), és
kiszűri az érvénytelen koordinátájú markereket, illetve a három pontnál kevesebbel rendelkező
sokszögeket -- így a hibás adat nem törheti el a térképet. A teljes kliens--szerver folyamatot
a~\ref{fig:map-flow}.~ábra foglalja össze.

% TODO: M-01 folyamatábra (figures/m-01-map-flow.pdf) — rétegválasztás → scope-kérés →
%       tagság-ellenőrzés → markerek/polygonok → rajzolás → mentés Seq-pontokkal.
\begin{figure}[htbp]
	\centering
	\fbox{\parbox[c][6.5cm][c]{0.9\textwidth}{\centering M-01: A térképi adatkezelés folyamatábrája (elkészítendő)}}
	\caption{A térképi adatkezelés folyamata a rétegválasztástól a mentésig}
	\label{fig:map-flow}
\end{figure}

\section{Adatfolyam, biztonság és korlátok}
\label{sec:map-korlatok}

A hozzáférés-szabályozás rétegei jól elválnak: a hatókör-alapú olvasás a személyes adatokat a
létrehozóra, a társasági réteget a tagságra korlátozza; az írás bejelentkezéshez és társasági
objektumnál tagsághoz kötött; a módosítás és törlés pedig szigorúan a létrehozóé. Ebből két
tudatosan vállalt, de dokumentálandó korlát is következik. Egyrészt a \emph{láthatósági kód
jelenleg címkeként működik}: a \texttt{PUBLIC}/\texttt{MEMBERS}/\texttt{ADMIN} érték eltárolódik
és a felületen megjelenik, a lekérdezések azonban nem szűrnek rá -- a tényleges elérhetőséget
kizárólag a hatókör-szabályok döntik el. Másrészt a \texttt{CanModify} szabály az adminisztrátort
sem engedi kivételezni, és a társaság tulajdonosa sem kezelheti tagjai jelöléseit; a szolgáltatás
kódjában elhelyezett megjegyzés szerint a szerepalapú finomítás későbbi bővítési pont.

További korlátok: a kliens a meglévő objektumok \emph{szerkesztését} nem kínálja fel -- a
módosító végpontok és a kliensoldali API-függvények (\texttt{updateMarker},
\texttt{updatePolygon}) készen állnak, felületi művelet azonban csak a létrehozás és a törlés --;
a térképoldalhoz végponttól végpontig tartó teszt nem készült, mert a Google Maps betöltése
API-kulcshoz kötött; az útvonalak (\texttt{MapPath}) pedig a~\ref{sec:map-adatmodell}.~alfejezetben
leírtak szerint csak előkészítettek. E pontok a továbbfejlesztési fejezetben
(\ref{chap:osszefoglalas}.~fejezet) térnek vissza.

% !TeX root = ../main.tex
% 10. fejezet — AI-chatbot és időjárás-integráció
% Hivatkozott források: Backend/HUNting.API/Controllers/ChatController.cs,
%   Services/OpenAiService.cs, Models/Chat/ChatRequest.cs, Program.cs,
%   Frontend/HUNting/src/components/Chatbot.tsx, services/weatherService.ts,
%   pages/Weather/WeatherPage.tsx

\chapter{AI-chatbot és időjárás-integráció}
\label{chap:chatbot}

A gerinc utolsó fejezete a két külső szolgáltatásra épülő kiegészítő funkciót tárgyalja: a
mesterséges intelligenciára épülő vadászati asszisztenst és az időjárás-lekérdezést. A fejezet
egyben a funkciófejezetek egységes sablonjának határesete is: e funkciók \emph{nem használnak
saját adatbázis-entitást} -- a chatbot állapotmentes (a beszélgetés csak a kliens memóriájában
él), az időjárás pedig teljes egészében kliensoldali integráció --, ezért az adatmodell-szakasz
helyét ennek indoklása veszi át, a hangsúly pedig a külső API-integrációk mérnöki részleteire
kerül.

\section{A chatbot célja és felhasználói működése}
\label{sec:chat-cel}

A „Vadászati Asszisztens” az alkalmazás minden fő oldalán elérhető, lebegő gombból nyitható
csevegőablak, amely magyar nyelven válaszol vadászati, vadgazdálkodási és fegyverismereti
kérdésekre -- tematikusan szűkítve az 1996. évi LV. törvény (a vad védelméről, a
vadgazdálkodásról, valamint a vadászatról) és a fegyvertartási szabályozás körére, kiegészítve
vadreceptekkel és időjárási témákkal. A felhasználó kötetlen szöveges kérdést tesz fel, a válasz
formázott (címsorokat, felsorolásokat is tartalmazó) szövegként jelenik meg; a beszélgetés a
munkamenet alatt folytatólagos, az ablak bezárásáig a korábbi üzenetek is láthatók. A widget a
hitelesítési oldalakon nem jelenik meg, mert azokat külön elrendezés (\texttt{AuthLayout})
szolgálja ki, amely nem tartalmazza a komponenst.

% TODO: SC-15 képernyőkép (figures/sc-15-chatbot.png) — nyitott widget magyar válasszal; OpenAI-kulcs kell.
\begin{figure}[htbp]
	\centering
	\fbox{\parbox[c][5.5cm][c]{0.85\textwidth}{\centering SC-15: A chatbot-widget megnyitva (elkészítendő)}}
	\caption{A Vadászati Asszisztens csevegőablaka}
	\label{fig:chatbot-screen}
\end{figure}

\section{A chatbot backend-megvalósítása}
\label{sec:chat-backend}

A \texttt{ChatController} egyetlen végpontot kínál (\texttt{POST /api/Chat}), amely a
\texttt{ChatRequest} DTO-t fogadja: az új üzenetet és az eddigi beszélgetés üzenetlistáját
(szerep--tartalom párok). Üres üzenetre 400-as hibával válaszol, az integráció bármely hibáját
pedig általános, magyar nyelvű 500-as üzenetté alakítja -- a belső hibarészletek így nem
szivárognak ki a kliens felé. Az érdemi munkát az \texttt{OpenAiService} végzi, amely a
\texttt{Program.cs}-ben típusos HTTP-kliensként van regisztrálva
(\texttt{AddHttpClient<IOpenAiService, OpenAiService>}), így a \texttt{HttpClient} életciklusát a
keretrendszer gyára kezeli.

A szolgáltatás működése négy lépésre bontható. Először a konfigurációból kiolvassa az
\texttt{OpenAI:ApiKey} kulcsot; ennek hiányában naplózott kivétellel leáll -- a kulcs tehát nem a
kódban, hanem a konfigurációban (fejlesztői gépen user secrets / \texttt{appsettings}) él.
Másodszor összeállítja az üzenetlistát: az élére egy részletes magyar \emph{rendszerprompt} kerül,
amely kijelöli az asszisztens szakterületét, megköveteli a témán kívüli kérdések udvarias
elutasítását, és kifejezetten tiltja a nem létező jogszabályi adatok kitalálását. Harmadszor a
kliens által küldött előzményeket \emph{szanitálva} fűzi be: a szerepmezőt fehérlistával
ellenőrzi, és minden nem \texttt{user} vagy \texttt{assistant} értéket \texttt{user} szerepre ír
át (\ref{lst:openai-request}.~kódrészlet) -- e nélkül egy rosszindulatú kliens
\texttt{system} szerepű üzenettel felülbírálhatná a rendszerpromptot (prompt-injekció). Végül a
kérés a \texttt{gpt-4o} modellel, alacsony (0{,}2) hőmérséklettel -- azaz konzervatív,
tényközpontú beállítással -- és 1500 tokenes válaszkorláttal megy az OpenAI
chat-completions végpontjára, Bearer-fejlécben vitt kulccsal.

\begin{lstlisting}[language={[Sharp]C}, caption={Az előzmények szanitálása és a kérés összeállítása az \texttt{OpenAiService} osztályban (kivonat, megjegyzések nélkül)}, label={lst:openai-request}]
var messages = new List<object>
{
    new { role = "system", content = systemPrompt }
};

foreach (var msg in history)
{
    var role = msg.Role == "user" || msg.Role == "assistant" ? msg.Role : "user";
    messages.Add(new { role = role, content = msg.Content });
}

messages.Add(new { role = "user", content = newMessage });

var requestBody = new
{
    model = "gpt-4o",
    messages = messages,
    temperature = 0.2,
    max_tokens = 1500
};

using var request = new HttpRequestMessage(HttpMethod.Post,
    "https://api.openai.com/v1/chat/completions");
request.Headers.Authorization = new AuthenticationHeaderValue("Bearer", apiKey);
request.Content = new StringContent(JsonSerializer.Serialize(requestBody),
    Encoding.UTF8, "application/json");
\end{lstlisting}

A válaszfeldolgozás \texttt{JsonDocument}-tel navigál a
\texttt{choices[0].message.content} útvonalon; sikertelen HTTP-státusz esetén a szolgáltatás a
hibatörzset naplózza, majd általános kivételt dob, amelyet a kontroller fordít 500-as válaszra. A
beszélgetés a szerveren \emph{nem perzisztálódik}: a \texttt{ChatRequest} és a
\texttt{ChatMessage} tisztán átviteli objektum, adatbázistábla nem tartozik hozzájuk -- minden
kérés önmagában hordozza a kontextusát. A kérés teljes életútját
a~\ref{fig:chat-sequence}.~ábra mutatja.

% TODO: C-01 szekvenciadiagram (figures/c-01-chat-sequence.pdf) — Chatbot → POST /Chat →
%       OpenAiService → api.openai.com → reply; szerep-szanitálás kiemelve.
\begin{figure}[htbp]
	\centering
	\fbox{\parbox[c][6cm][c]{0.9\textwidth}{\centering C-01: A chatbot-kérés feldolgozásának szekvenciadiagramja (elkészítendő)}}
	\caption{A chatbot-kérés útja a widgettől az OpenAI API-ig és vissza}
	\label{fig:chat-sequence}
\end{figure}

\section{A chatbot frontend-megvalósítása}
\label{sec:chat-frontend}

A \texttt{Chatbot} komponens lebegő gombként ül az elrendezés bal alsó sarkában, és kattintásra
csevegőablakká nyílik. A beszélgetést egyszerű komponens-állapot tárolja: az üzenetlista egy
asszisztensi üdvözlő üzenettel indul, a beküldött kérdés azonnal bekerül a listába, a
betöltésjelző pedig a válasz megérkezéséig forog. A kérés a közös \texttt{apiClient} példányon
megy ki, az előzménylistából a nyitó üdvözlést levágva (\texttt{history: chatHistory.slice(1)}),
mert az nem a modelltől származik. Hiba esetén a widget általános hibaszöveget fűz a
beszélgetésbe, így a felhasználói élmény nem törik meg. A válaszok \texttt{react-markdown}
komponenssel renderelődnek, egyedi elemleképezéssel (címsorok, listák, kiemelések) -- a modell
által adott strukturált válaszok így olvashatóan jelennek meg. Két ergonómiai részlet teszi
teljessé: az üzenetlista minden változáskor az aljára gördül, a küldés pedig Enterrel is
kiváltható, üres bevitelnél letiltva. A beszélgetés nem kerül sem \texttt{localStorage}-ba, sem
szerverre: az oldal újratöltése új beszélgetést kezd.

\section{Időjárás-integráció}
\label{sec:weather}

Az időjárás-oldal a vadászat tervezését segíti: városnév alapján mutatja az aktuális
hőmérsékletet, hőérzetet, páratartalmat, a szöveges leírást (magyar nyelven) és a szélsebességet
-- utóbbit a szélirányt jelző, a \texttt{wind\_deg} értékkel elforgatott nyíllal. A megvalósítás
a többi funkciótól eltérő, tudatos architektúra-kompromisszum: a \texttt{weatherService} modul
\emph{közvetlenül}, a backend megkerülésével hívja az OpenWeather „current weather” végpontját
(\texttt{q}, \texttt{units=metric}, \texttt{lang=hu} paraméterekkel), saját axios-hívással, nem a
közös \texttt{apiClient}-tel -- hiszen itt sem hitelesítési fejlécre, sem a backend 401-kezelésére
nincs szükség. Az API-kulcsot a \texttt{VITE\_OPEN\_WEATHER\_API} környezeti változó adja. Az
oldal állapotkezelése szándékosan egyszerű: város, betöltés, hiba és eredmény lokális
állapotban, React Query nélkül -- egyszeri, gombnyomásra induló lekérdezésről lévén szó a
gyorsítótárazás nem adna hozzáadott értéket.

% TODO: SC-11 képernyőkép (figures/sc-11-weather.png) — találattal; OpenWeather-kulcs kell.
\begin{figure}[htbp]
	\centering
	\fbox{\parbox[c][5cm][c]{0.85\textwidth}{\centering SC-11: Az időjárás-oldal találattal (elkészítendő)}}
	\caption{Az időjárás-oldal egy lekérdezett város aktuális adataival}
	\label{fig:weather-screen}
\end{figure}

\section{Biztonsági megfontolások és korlátok}
\label{sec:chat-korlatok}

A chatbot legfontosabb védelmi vonala a szerveroldali elhelyezés: az OpenAI-kulcs soha nem jut el
a kliensig, a rendszerprompt nem módosítható kívülről, a szerep-fehérlista pedig a
prompt-injekció alapvető vektorát zárja le. Ugyanakkor a végponton \emph{nincs autorizáció és
nincs sebességkorlátozás}: anonim hívó is költséget generálhat a fizetős OpenAI-kvóta terhére --
éles üzem előtt bejelentkezéshez kötés és kérésszám-korlát (rate limiting) szükséges. Az
időjárás-integráció tükörképe ennek a kompromisszumnak: a kliensoldali hívás egyszerű és nem
terheli a backendet, cserébe az OpenWeather-kulcs a böngészőbe kerül, azaz kiolvasható -- ingyenes
kvótájú kulcsnál vállalható, éles környezetben backend-proxy javasolt. További apró, ismert
hibák: hiányzó kulcs esetén a szolgáltatás csendben \texttt{null} értékkel tér vissza, így a
felület hibajelzés nélkül marad üres, az időjárás-ikon pedig rögzített \texttt{http://} sémájú
URL-ről töltődik, amelyet a böngésző HTTPS alatt vegyes tartalomként blokkolhat. E pontok a
továbbfejlesztési irányok (\ref{chap:osszefoglalas}.~fejezet) listáját gazdagítják.

% !TeX root = ../main.tex
% 11. fejezet — Tesztelés és üzemeltetés
% A ftsrg/thesis-template-latex sablonba \include-dal illesztendő tartalomfájl.
% Hivatkozott források: Backend/HUNting.API.Tests/* (HUNting.API.Tests.csproj,
%   Controllers/*Tests.cs, Services/OpenAiServiceTests.cs),
%   Frontend/HUNting/vitest.config.ts, src/setupTests.ts, src/**/*.test.tsx,
%   playwright.config.ts, e2e/*.spec.ts, e2e/helpers.ts,
%   Backend/HUNting.API/Properties/launchSettings.json, appsettings*.json, HUNting.API.csproj

\chapter{Tesztelés és üzemeltetés}
\label{chap:teszteles}

A rendszerhez háromszintű automatizált tesztkészlet tartozik, amely a tesztpiramis klasszikus
elrendezését követi (\ref{fig:tesztpiramis}.~ábra): az alapot 124 backend-egységteszt adja, a
középső szintet 17 kliensoldali komponenstesztfájl, a csúcsot pedig 26 végponttól végpontig
(end-to-end, E2E) tartó Playwright-teszt. A fejezet a három szint stratégiáját és jellegzetes
mintáit mutatja be, majd a helyi futtatás és konfiguráció gyakorlati kérdéseit, végül az
üzemeltetés és a tesztelés ismert korlátait összegzi. Fontos módszertani megjegyzés: a fejezet
a tesztek \emph{felépítését} dokumentálja a forráskód alapján; futási eredményeket
(sikerességi arányt, futásidőt) nem közöl.

% TODO: Q-01 ábra (figures/q-01-tesztpiramis.pdf) — piramis: 124 xUnit / 17 Vitest-tesztfájl /
% 26 Playwright E2E-teszt.
\begin{figure}[htbp]
	\centering
	\fbox{\parbox[c][5cm][c]{0.7\textwidth}{\centering Q-01: A HUNting tesztpiramisa
	(elkészítendő)}}
	\caption{A tesztkészlet megoszlása a tesztpiramis szintjei szerint}
	\label{fig:tesztpiramis}
\end{figure}

\section{Backend-egységtesztek}
\label{sec:teszt-backend}

A backend-teszteket az xUnit keretrendszer futtatja~\cite{xunitdocs}, a függőségek
helyettesítését a Moq könyvtár végzi; az adatelérés izolálására a tesztkészlet az EF Core
InMemory szolgáltatóját és a MockQueryable kiegészítőt használja. A 124 teszt kilenc
tesztosztályban oszlik el: a nyolc kontrollernek egy-egy osztály felel meg (a legnagyobbak a
társasági, a térképi és a hirdetési kontroller tesztjei, 25, 24, illetve 23 esettel), a
kilencedik az \texttt{OpenAiService} szolgáltatást fedi le.

A tesztek a lefedett komponens függőségi szerkezetéhez igazodva három izolációs stratégiát
alkalmaznak. Ahol a kontroller szolgáltatásinterfészre támaszkodik (hirdetések, térkép, chat),
ott a teszt az interfészt mockolja (\texttt{Mock<IAdService>},~stb.), és kizárólag a
kontroller HTTP-szintű viselkedését -- státuszkódokat, jogosultsági ágakat -- ellenőrzi. Ahol a
kontroller közvetlenül az adatbázis-kontextust használja (hitelesítés, hírek, társaságok,
vadásznapló), ott a teszt EF Core InMemory adatbázist épít, és a teljes adatutat lefedi. A
felhasználó-adminisztrációs tesztek a MockQueryable segítségével az Identity
\texttt{UserManager} lekérdezhető felületét helyettesítik. Mindhárom stratégiában közös a
hitelesített hívó szimulálása: a tesztek kézzel összeállított \texttt{ClaimsPrincipal}
objektumot (felhasználó-azonosító és szerep claimekkel) helyeznek a kontroller
\texttt{HttpContext}-jébe, így a jogosultsági ágak -- tulajdonos, idegen felhasználó,
adminisztrátor -- külön-külön tesztelhetők.

Külön említést érdemel a külső HTTP-függőség tesztelése. Az \texttt{OpenAiService} típusos
\texttt{HttpClient}-tel dolgozik, ezért a tesztek nem a klienst, hanem az az alatt ülő
\texttt{HttpMessageHandler} absztrakt osztályt mockolják
(\ref{lst:httphandler-mock}.~kódrészlet): így a szolgáltatás valódi \texttt{HttpClient}-tel fut,
miközben a hálózati hívás helyett előre összeállított válasz érkezik, a kimenő kérés pedig
elfogható és vizsgálható.

\begin{lstlisting}[language={[Sharp]C}, caption={A \texttt{HttpMessageHandler} mockolása az \texttt{OpenAiServiceTests}-ben (kivonat)}, label={lst:httphandler-mock}]
var handlerMock = new Mock<HttpMessageHandler>();
handlerMock.Protected()
    .Setup<Task<HttpResponseMessage>>(
        "SendAsync",
        ItExpr.IsAny<HttpRequestMessage>(),
        ItExpr.IsAny<CancellationToken>())
    .Callback<HttpRequestMessage, CancellationToken>((req, _) =>
    {
        captureRequest?.Invoke(req);
        // a keres torzse meg a dispose elott kerul kiolvasasra
    })
    .ReturnsAsync(response);

return new OpenAiService(new HttpClient(handlerMock.Object),
                         configuration,
                         NullLogger<OpenAiService>.Instance);
\end{lstlisting}

Erre a mintára épül a tesztkészlet egyik legértékesebb esete: a kérésnapló
szerep-szanitálásának ellenőrzése. A teszt \texttt{"hacker"} szerepű üzenetet csempész az
előzménylistába, majd az elfogott kéréstörzs JSON-jából bizonyítja, hogy a szolgáltatás azt
\texttt{"user"} szerepre írta át -- vagyis a kliens nem tud rendszerszintű utasítást
befecskendezni az OpenAI-kérésbe (\ref{chap:chatbot}.~fejezet).

\section{Frontend-komponenstesztek}
\label{sec:teszt-frontend}

A kliensoldali teszteket a Vitest futtatja~\cite{vitestdocs}, a komponensek viselkedésalapú
vizsgálatát a Testing Library végzi~\cite{testinglibrary}, a böngészőkörnyezetet pedig a
könnyűsúlyú happy-dom szimulálja; a konfiguráció (\texttt{vitest.config.ts}) a Vite
beállításaiból származik, így az útvonal-aliasok a tesztekben is működnek. A 17 tesztfájl a
komponensfa minden szintjét érinti: az alkalmazásgyökértől (\texttt{App.test.tsx}) az
elrendezéselemeken és UI-primitíveken át az oldalkomponensekig, amelyek a készlet zömét adják
(hitelesítési, hirdetési, társasági, térképi, napló-, időjárás- és profiloldalak).

A tesztek közös mintája a határfelületek modulszintű mockolása: a \texttt{vi.mock} hívás az
API-szolgáltatásmodulokat (például \texttt{adsApi}) és a \texttt{useAuth} hookot helyettesíti,
így a komponens valódi hálózati hívás és valódi hitelesítési állapot nélkül, előre definiált
adatokkal renderelhető. Minden teszt friss \texttt{QueryClient} példányt kap kikapcsolt
újrapróbálkozással, a renderelés pedig a valódi provider-fát (QueryClientProvider, Router)
reprodukálja -- a tesztek tehát nem implementációs részleteket, hanem a felhasználó által
látható viselkedést (megjelenő szövegeket, gombokra adott reakciókat) ellenőrzik. A térképoldal
tesztje a Google Maps betöltését is mockolja (\texttt{useLoadScript}), így API-kulcs nélkül is
futtatható.

\section{Végponttól végpontig tesztek}
\label{sec:teszt-e2e}

A legfelső szintet a Playwright-tesztek adják~\cite{playwrightdocs}: hét tesztfájl (smoke,
hitelesítés, hirdetések, társaságok, vadásznapló, chatbot, admin) összesen 26 tesztesete
valódi böngészőben (Chromium), a teljes, integrált rendszer -- frontend, backend és adatbázis
-- ellen fut. A konfiguráció kulcseleme a \texttt{webServer} blokk
(\ref{lst:webserver}.~kódrészlet): a tesztfuttató maga indítja el a backendet és a frontendet,
ha még nem futnak, így a készlet egyetlen paranccsal, külön előkészítés nélkül futtatható.

\begin{lstlisting}[caption={A Playwright \texttt{webServer} konfigurációja (kivonat)}, label={lst:webserver}]
webServer: [
  {
    command: "dotnet run --project ../../Backend/HUNting.API" +
             " --launch-profile http",
    url: "http://localhost:5137/api/ads",
    reuseExistingServer: true,
    timeout: 180_000,
  },
  {
    command: "npm run dev -- --port 5173 --strictPort",
    url: "http://localhost:5173",
    reuseExistingServer: true,
    env: {
      // az apiClient.ts dev fallbackje nem a valos backend port
      VITE_API_BASE_URL: "https://localhost:7003/api",
    },
  },
],
\end{lstlisting}

A konfiguráció több, a rendszer sajátosságaiból fakadó döntést rögzít. A tesztek közös
adatbázison dolgoznak, ezért a párhuzamosság kikapcsolt (\texttt{workers:~1}), a futás
szekvenciális. A backend önaláírt fejlesztői tanúsítványa miatt a HTTPS-hibák elfogadottak
(\texttt{ignoreHTTPSErrors}). A frontend indításakor a konfiguráció a
\texttt{VITE\_API\_BASE\_URL} változót a backend tényleges 7003-as portjára állítja, ezzel
megkerülve az \texttt{apiClient} téves fejlesztői alapértelmezését
(\ref{sec:alapok-frontend}.~alfejezet).

A tesztadat-kezelés kétféle fiókra épül. A csak olvasó forgatókönyvek a kezdeti adatfeltöltés
ismert fiókjait (adminisztrátor, seedelt vadász) használják; a bejelentkezés nem a felületen,
hanem API-híváson át történik, és a kapott tokent a segédmodul a frontend
\texttt{localStorage}-kulcsaiba injektálja -- így a tesztek gyorsabbak, és nem függenek a
bejelentkező űrlaptól. Az állapotot módosító, csak egyszer végrehajtható forgatókönyvekhez
(például első tagfelvételi kérelem) a segédmodul minden futásban friss, időbélyeg-alapú egyedi
e-mail-címmel regisztrált fiókot hoz létre, ezzel biztosítva az ismételhetőséget. A térképoldal
E2E-tesztje tudatosan hiányzik: a Google Maps betöltése API-kulcshoz kötött
(\ref{chap:terkep}.~fejezet).

\section{Helyi futtatás és konfiguráció}
\label{sec:teszt-uzemeltetes}

A rendszer éles üzemeltetésre nincs felkészítve; a fejlesztői és demonstrációs futtatás
lépései és feltételei a következők. A backend a \texttt{dotnet run} paranccsal indul, a
fejlesztői profil szerint a \texttt{https://localhost:7003} és a \texttt{http://localhost:5137}
címeken (\texttt{launchSettings.json}); induláskor automatikusan érvényesíti a migrációkat és
lefuttatja az adatfeltöltést (\ref{sec:alapok-adat}.~alfejezet), fejlesztői módban pedig a
Swagger-felület is elérhető. Az adatbázishoz SQL Server LocalDB szükséges; a HTTPS-hez a .NET
fejlesztői tanúsítványa (\texttt{dotnet dev-certs https --trust}). A frontend a Vite fejlesztői
kiszolgálóján, az 5173-as porton fut (\texttt{npm run dev}); ezt a származási helyet engedélyezi
a backend CORS-házirendje.

A külső szolgáltatásokhoz három kulcs szükséges, három különböző helyen: az OpenAI-kulcs a
backend konfigurációjában (\texttt{OpenAI:ApiKey} -- fejlesztői gépen a user secrets tárolóban
vagy az \texttt{appsettings.Development.json}-ban), a Google Maps- és az OpenWeather-kulcs a
frontend környezeti változóiban (\texttt{VITE\_GOOGLE\_MAPS\_API\_KEY}, illetve
\texttt{VITE\_OPEN\_WEATHER\_API}). Kulcs hiányában a rendszer többi része működőképes marad:
a chatbot hibaüzenetet ad, a térképoldal betöltési hibát jelez, az időjárásoldal üres marad.

\section{Ismert korlátok}
\label{sec:teszt-korlatok}

A tesztelés és az üzemeltetés őszinte állapotfelmérése a következő pontokat rögzíti.
Folyamatos integrációs (CI) futtatókörnyezet nincs: a tesztek kézi indításúak, és a dolgozat --
a fenti módszertani megjegyzésnek megfelelően -- futási eredményeket nem dokumentál.
A teszt- és a fő projekt célkeretrendszere eltér (a tesztprojekt .NET~9-re, az API .NET~8-ra
fordul), ami a tesztek futtathatóságát nem akadályozza, de egységesítése indokolt. Éles
telepítési folyamat, konténerizáció és környezetenkénti konfigurációkezelés nem készült; az
\texttt{appsettings.json}-ban fejlesztői helykitöltő JWT-kulcs szerepel. A biztonsági
adósságok~\cite{owasp2021} közül kiemelendő a tokenvisszavonás és a frissítő token hiánya
(\ref{chap:felhasznalok}.~fejezet), valamint az anonim chatvégpont és a kliensoldali
API-kulcsok kitettsége (\ref{chap:chatbot}.~fejezet). Ezek a pontok a továbbfejlesztési
irányok között (\ref{chap:osszefoglalas}.~fejezet) térnek vissza.

% !TeX root = ../main.tex
% 12. fejezet — Összefoglalás és továbbfejlesztési lehetőségek
% A ftsrg/thesis-template-latex sablonba \include-dal illesztendő tartalomfájl.
% A fejezetre a 7., 8., 9. és 10. fejezet hivatkozik (chap:osszefoglalas) — a címke KÖTELEZŐ.
% Hivatkozott források: a dolgozat 5–11. fejezetei és a bennük dokumentált korlátok.

\chapter{Összefoglalás és továbbfejlesztési lehetőségek}
\label{chap:osszefoglalas}

\section{Eredmények}
\label{sec:ossz-eredmenyek}

A dolgozat egy teljes, működő webalkalmazás -- a HUNting vadászati közösségi és hirdetési
portál -- megtervezését és megvalósítását mutatta be, az adatbázistól a felhasználói felületig.
A bevezetésben kitűzött hat funkcióterület mindegyike megvalósult, a gerincfejezetekben
dokumentált terjedelemben. A \emph{felhasználókezelés} a kettős fióktípusra (egyéni vadász és
vadásztársaság) épülő regisztrációt, JWT-alapú hitelesítést, két szintű -- szerepkör- és
erőforrástulajdonos-alapú -- jogosultsági modellt és admin-felületet ad
(\ref{chap:felhasznalok}.~fejezet). A \emph{hirdetési piactér} kategorizált, szűrhető és
lapozott hirdetéslistát, képfeltöltéses hirdetéskezelést és kettős (felhasználói vagy
társasági) tulajdonlást valósít meg (\ref{chap:hirdetesek}.~fejezet). A \emph{társasági modul}
a tagfelvételi kérelmek teljes életciklusát, társaságon belüli szerepeket és tagoknak szóló
bejegyzéseket kezel (\ref{chap:tarsasagok}.~fejezet). A \emph{vadásznapló} képmelléklettel
rögzíti a személyes vadászati eseményeket (\ref{chap:naplo}.~fejezet). A \emph{térképes modul}
Google Maps-alapú megjelenítéssel, személyes és társasági rétegeken kezel jelölőket és
rajzolt területeket (\ref{chap:terkep}.~fejezet). Végül az \emph{AI-chatbot} szerveroldali
OpenAI-integrációval, szerep-szanitálással ad vadászati asszisztenst, az
\emph{időjárásoldal} pedig aktuális időjárási adatokat jelenít meg
(\ref{chap:chatbot}.~fejezet).

A megvalósítás mérnöki értékét a közös alapok adják: az egységes, 26 táblás, kifejezetten
megtervezett törlési viselkedésekkel és megszorításokkal bíró adatmodell, az idempotens
kezdeti adatfeltöltés, a következetes middleware-csővezeték és a kliensoldali közös
HTTP-réteg (\ref{chap:alapok}.~fejezet). A kódbázishoz háromszintű tesztkészlet készült --
124 backend-egységteszt, 17 komponenstesztfájl és 26 végponttól végpontig teszt --, amely a
jogosultsági ágakat, a külső integrációk izolálását és a kritikus felhasználói utakat fedi le
(\ref{chap:teszteles}.~fejezet). A dolgozat végig törekedett az őszinte minősítésre: minden
funkciófejezet külön szakaszban rögzítette a megvalósítás korlátait, megkülönböztetve a
működő, az előkészített és a hiányzó képességeket.

\section{Továbbfejlesztési lehetőségek}
\label{sec:ossz-tovabbfejlesztes}

A funkciófejezetekben dokumentált korlátok egyben a továbbfejlesztés térképét is kirajzolják.
A legfontosabb irányok a következők.

\begin{itemize}
	\item \emph{Hírfolyam befejezése:} a backend-oldalon előkészített hírmodul
	      (\ref{chap:tarsasagok}.~fejezet) felhasználói felületének elkészítése, a hiányzó
	      autorizáció pótlásával együtt -- ez a rendszer legégetőbb biztonsági adóssága.
	\item \emph{Hitelesítés megerősítése:} frissítő token és tokenvisszavonás bevezetése, a
	      jelszószabályzat szigorítása, a konfigurációs titkok (JWT-kulcs, API-kulcsok)
	      titokkezelőbe szervezése.
	\item \emph{Térképi funkciók kiteljesítése:} az adatmodell szintjén előkészített útvonalak
	      (\texttt{MapPath}) végpontjainak és felületének elkészítése; a láthatósági kódok
	      tényleges szűrési logikájának megvalósítása; szerkesztőfelület a meglévő
	      PUT-végpontokhoz (\ref{chap:terkep}.~fejezet).
	\item \emph{Felületi hiányok pótlása:} szerkesztő- és törlőműveletek a vadásznaplóban,
	      lapozás a naplólistában (\ref{chap:naplo}.~fejezet), admin-felület a létrehozó és
	      módosító végpontokhoz, valamint a felületi szövegek ékezeteinek rendezése.
	\item \emph{Külső integrációk keményítése:} az anonim chatvégpont sebességkorlátozása és a
	      kliensoldali OpenWeather-hívás backend-proxy mögé szervezése
	      (\ref{chap:chatbot}.~fejezet).
	\item \emph{Közösségi funkciók bővítése:} értesítések (például tagfelvételi döntésről),
	      tagok közötti üzenetküldés, valamint hosszabb távon natív mobilalkalmazás.
	\item \emph{Üzemeltetési érettség:} CI-futtatás a meglévő tesztkészletre, konténerizálás,
	      környezetenkénti konfiguráció és telepítési folyamat; kódszinten a lusta
	      oldalbetöltés bekapcsolása és a kliensoldali típusok generálása az
	      OpenAPI-leírásból.
\end{itemize}

A felsorolás terjedelme nem a rendszer hiányosságát, hanem a választott megközelítés
következményét tükrözi: a dolgozat egy szélesen integrált rendszert valósított meg hat
funkcióterületen, teljes vertikumban, és minden területen pontosan dokumentálta, hol húzódik a
megvalósítás határa. Az elkészült alap -- a közös adatmodell, a jogosultsági rendszer és a
tesztkészlet -- úgy lett kialakítva, hogy ezek az irányok töréspont nélkül építhetők rá.


\bibliographystyle{plain}
\bibliography{references}

\end{document}
